\documentclass{article}

%\usepackage{sectsty}
\usepackage{titlesec}
\usepackage{url}

\usepackage{eso-pic}

\usepackage{setspace}
%\doublespacing
%\onehalfspacing


% line
\usepackage[object=vectorian]{pgfornament} %%  http://altermundus.com/pages/tkz/ornament/index.html

\definecolor{darkRed}{rgb}{.2,.0,.1}

\newcommand{\decoline}{%

\vspace{-2em}
{\color{darkRed!60!cyan}\noindent\hfil{\EnglischeLinie}\hfil}
%{\color{darkRed!60!cyan}\noindent\hfil\rule{0.5\textwidth}{.8pt}\hfil}
\vspace{-2.25em}}


\newcommand{\whdecoline}{%

\vspace{-2em}
{\color{white}\noindent\hfil{\whEnglischeLinie}\hfil}
%{\color{darkRed!60!cyan}\noindent\hfil\rule{0.5\textwidth}{.8pt}\hfil}
\vspace{-2.25em}}

\newcommand{\sectionline}[1]{%
  \noindent
  \begin{center}
  {\color{#1}
    \resizebox{0.5\linewidth}{1ex}
    {{%
    {\begin{tikzpicture}
    \node  (C) at (0,0) {};
    \node (D) at (9,0) {};
    \path (C) to [ornament=84] (D);
    \end{tikzpicture}}}}}%
    \end{center}
  }
  
\newcommand{\EnglischeLinie}{
\sectionline{darkRed!60!cyan}
}

\newcommand{\whEnglischeLinie}{
\sectionline{white}
}
% ///


\newcommand{\sectionlinenoc}[2]{%
	\noindent
		{\color{#1}
			\resizebox{#2}{1ex}
			{{%
					{\begin{tikzpicture}
						\node  (C) at (0,0) {};
						\node (D) at (9,0) {};
						\path (C) to [ornament=84] (D);
						\end{tikzpicture}}}}}%
}

\newcommand{\shortsectionline}[1]{%
	\noindent
	\begin{center}
		{\color{#1}
			\resizebox{.2\linewidth}{1.5ex}
			{{%
					{\begin{tikzpicture}
						\node  (C) at (0,0) {};
						\node (D) at (9,0) {};
						\path (C) to [ornament=84] (D);
						\end{tikzpicture}}}}}%
	\end{center}
}

\definecolor{blGreen}{rgb}{.2,.7,.3}

\newcommand{\shortdecoline}{\vspace*{-.65em}\shortsectionline{blGreen!10!orange}\vspace*{-.85em}}
\newcommand{\shortdecolineadj}[2]{\vspace*{#1}\shortsectionline{blGreen!10!orange}\vspace*{#2}}


\usepackage[flushmargin]{footmisc}

\usepackage[letterpaper, left=.45in,right=.45in,top=1in,bottom=1in]{geometry}

\colorlet{codegr}{black!80!blue}

\setlength{\columnsep}{7mm}

%% rightx
\newcommand{\htwoo}{H$_2$O}
\newcommand{\dtwoo}{D$_2$O}
\newcommand{\POne}{$P_1$}
\newcommand{\PTwo}{$P_2$}
\newcommand{\Pone}{$P_1$}
\newcommand{\Ptwo}{$P_2$}
\newcommand{\Pprop}{$P$}

\newcommand{\Stwo}{$S_2$}
\newcommand{\Sone}{$S_1$}

\newcommand{\facondeparler}{\textit{fa{\c{c}}on de parler}}

\newcounter{sentenceCounter}{}

\newcommand{\writeIncSentenceCounter}{\refstepcounter{sentenceCounter}(\arabic{sentenceCounter})}{}


\newcommand{\listItemMark}{\rotatebox{30}{\raisebox{2pt}{\color{green!30!yellow!40!black}{\begin{tiny}$\blacktriangleright$\end{tiny}}}}}
\newenvironment{sentenceList}{\begin{list}{\listItemMark}{\setlength{\leftmargin}{.5em}\setlength{\itemsep}{-.1em}\setlength{\topsep}{.85em}}\begin{small}}{\end{small}\end{list}}

\newcommand{\sentenceItem}{\item \writeIncSentenceCounter}{}
% ///

\usepackage{etoolbox}

\AtBeginEnvironment{thebibliography}{\linespread{1}\selectfont}

%?\usepackage{mathptmx}


%?\titleformat*{\subsection}{\small\bfseries}

%?\usepackage{eufrak}
\usepackage{wasysym}
\usepackage{textcomp}
\usepackage{amssymb}

\usepackage{microtype}



\DeclareMathAlphabet{\mathcal}{OMS}{cmsy}{m}{n}
%\usepackage{euler}

\let\OldI\i


\newcommand{\mdash}{---}
\newcommand{\q}[1]{``#1"}
\newcommand{\sq}[1]{`#1'}
\renewcommand{\i}[1]{\textit{#1}}

\newcommand{\T}[1]{\raisebox{-2pt}{\ensuremath{\mathcal{T}}}\textit{\tiny #1}}
%\newcommand{\T}{\ensuremath{\mathscr{T2}}}

\newcommand{\TSupT}{\ensuremath{{\T2}\makebox[4pt][r]{\raisebox{5pt}{{\scalebox{.6}{\T1}}}}}}

\let\OldFootnoteSize\footnotesize
\renewcommand{\footnotesize}{\scriptsize}
	
%\newcommand{\Tnoindex}{\raisebox{-2pt}{\ensuremath{\mathfrak{t}}}}
\newcommand{\Tnoindex}{\ensuremath{\mathfrak{t}}}

\newcommand{\typeAbove}{%
\raisebox{-1pt}{\rotatebox{90}{\begin{tiny}$\diagdown$\makebox[1pt][c]{$\diagup$}\end{tiny}}}}

%\newcommand{\typeT}{\ensuremath{type\raisebox{.5pt}{\makebox[3pt][c]{-}}T}}
%\newcommand{\typeT}{\ensuremath{\mathcal{T}}}
\newcommand{\typeT}{\ensuremath{\mathfrak{t}}}

\newcommand{\TValues}{\typeT{}-values} 

\newcommand{\emigres}{\'emigr\'es}

\newcommand{\Retore}{Retor\'e}
\newcommand{\Aurelie}{Aur\'elie}
\newcommand{\Descles}{D\'escles}

\newcommand{\ala}{\`a la}

\newcommand{\picalculus}{\ensuremath{\pi}-calculus}

\newcommand{\qmarkdubious}{\raisebox{4pt}{{\footnotesize\textbf{?}}}}

%\newcommand{\TypeCat}{\ensuremath{\mathcal{T}}}
\newcommand{\TypeCat}{\ensuremath{\mathfrak{t}}}

\newcommand{\ADJplusNPeqNP}{\ensuremath{ADJ + NP = NP}}

%\newcommand{\outarrow}{\ensuremath{\overset{..}{\rightarrow}}}

\newcommand{\outarrow}{\makebox[-2pt][l]{%
\raisebox{4pt}{..}}\ensuremath{\rightarrow}}

\newcommand{\argarrow}{\hspace{1pt} \makebox[-2pt][l]{%
		\raisebox{4pt}{.}}\ensuremath{\rightarrow} \hspace{1pt}}


\newcommand{\smoutarrow}{\makebox[-1pt][l]{%
		\raisebox{3pt}{..}}\ensuremath{\rightarrow}}

\newcommand{\smargarrow}{\hspace{1pt} \makebox[-2pt][l]{%
		\raisebox{3pt}{.}}\ensuremath{\rightarrow} \hspace{1pt}}



%\newcommand{\NounToNoun}{N \outarrow N}
%\newcommand{\VisNtoS}{V :: N \outarrow N}

\newcommand{\mmbox}[1]{#1}

\newcommand{\NtoN}{\mmbox{\ensuremath{N} \outarrow{} \ensuremath{N}}}
\newcommand{\NounToNoun}{\mmbox{\ensuremath{N} \outarrow{} \ensuremath{N}}}
\newcommand{\VisNtoS}{\mmbox{\ensuremath{V :: N} \outarrow{} \Prop}}
\newcommand{\AdjisNtoN}{\mmbox{\hspace{2pt}\ensuremath{\mathcal{A}
\scalebox{.7}{\ensuremath{\mathcal{DJ}}} :: N} \outarrow{} \Prop}}


\newcommand{\thatPhrase}{\AcronymText{that}-phrase}
\newcommand{\thatPhrases}{\AcronymText{that}-phrases}

\usepackage{graphicx}

\newcommand{\rzauthor}[3]{{
{\vspace{-4em}}{\fontfamily{gar}\fontseries{b}\selectfont% 
{\begin{center}\textls*[103]{#1}\\ \vspace{-2em}%
{\scalebox{.7}{\fontfamily{phv}\fontshape{it}\selectfont\footnotesize #3}}#2\end{center}}

{\vspace{-2.25em}}
{\hfill\small\today}

{\vspace{.25em}}
}}}			 


\newcommand{\rzauth}[3]{{
		{\vspace{-1em}}{\fontfamily{gar}\fontseries{b}\selectfont% 
			{\begin{center}\textls*[103]{#1}\\ \vspace{1em}%
					{\scalebox{.7}{\fontfamily{phv}\fontshape{it}\selectfont\footnotesize #3}}#2\end{center}}
			
			{\vspace{-6em}}
			{\hfill\small {\raisebox{1em}{\today}}}
			
			{\vspace{0em}}
		}}}			 




\newcommand{\rztitle}[1]{\begin{center}\fontfamily{phv}\fontseries{b}\selectfont #1\end{center}}


\newcommand{\Jorgen}{J{\o}rgen}

\newcommand{\Prop}{\ensuremath{\mathcal{P}rop}}

\newcommand{\NN}{\ensuremath{N}\hspace{-1pt}\argarrow\hspace{-1pt}\ensuremath{N}}
\newcommand{\Npl}{%
\ensuremath{N}\raisebox{8pt}{\hspace{-1pt}\ensuremath{\rotatebox{180}{{\begin{footnotesize}$\dotplus$\end{footnotesize}}}}}}
\newcommand{\NPl}{\Npl}
\newcommand{\NtoNpl}{%
\mmbox{\ensuremath{N} \outarrow{} \Npl}}
\newcommand{\NpltoNpl}{\mmbox{\Npl{} \outarrow{} \Npl}}
\newcommand{\PropToN}{\mmbox{\Prop \hspace{1pt} \outarrow{} \N}}

\newcommand{\VisNNtoProp}{\mmbox{\ensuremath{V :: N} \argarrow \ensuremath{N} \outarrow{} \Prop}}
\newcommand{\VisNProptoProp}{\mmbox{\ensuremath{V :: N} \argarrow  \Prop{} \hspace{2pt} \outarrow{} \Prop}}


\newcommand{\NSing}{\ensuremath{N}\raisebox{5pt}{\scalebox{0.65}{\ensuremath{%
{\odot}}}}}


\newcommand{\Z}{\ensuremath{\mathbb{Z}}}

\newcommand{\VPpNPeqS}{\ensuremath{VP + NP = S}}
\newcommand{\NSingToNPl}{\mmbox{\NSing{} \outarrow{} \Npl}}

\newcommand{\NNtoS}{\mmbox{\NN{} \outarrow{} \Prop}}
\newcommand{\NNtoN}{\mmbox{\NN{} \outarrow{} \ensuremath{N}}}
\newcommand{\N}{\mmbox{\ensuremath{N}}}
\newcommand{\NtoProp}{\mmbox{\N{} \outarrow{} \Prop}}
\newcommand{\NNtoProp}{\mmbox{\NN{} \outarrow{} \Prop{}}}
\newcommand{\ProptoN}{\mmbox{\Prop{} \outarrow{} \N}}


\newcommand{\PropToNYieldsN}{\PropToN{} \raisebox{1pt}{\ensuremath{\textgreater\hspace{-5pt}\raisebox{.15pt}{\ensuremath{\Rightarrow}}}}{} \N}

\newcommand{\SeqNPplVP}{\mmbox{{\small\texttt{S = NP + VP}}}}

\newcommand{\Arg}{\ensuremath{A}\tiny{rg}}
\newcommand{\argsToReturn}{\mbox{\ensuremath{{\Arg}_0} \hspace{-4pt} \smargarrow 
		\hspace{-3pt}
\ensuremath{{\Arg}_1} 
\hspace{-4pt} \smargarrow \hspace{-2pt} \ensuremath{\cdot\cdot\cdot} \hspace{0pt} \smoutarrow  \hspace{1pt} \textit{Result}}}

\newcommand{\lowBlank}{\raisebox{-1pt}{\textthreequartersemdash\textthreequartersemdash\textthreequartersemdash}}

\newcommand{\BlankAfterBlank}{\lowBlank\texttt{after}\lowBlank}
\newcommand{\AfterNSingAndNSingToNPl}{\mbox{\texttt{after} :: %
\NSing \argarrow \NSing \hspace{1pt} \outarrow \hspace{1pt} \Npl}}

\newcommand{\archiveDate}[1]{{\footnotesize #1}}

\newcommand{\sentenceexample}[1]{
\begin{noindent}
\hspace*{-8pt}{{\raisebox{2pt}{{\scriptsize {\color{green!30!yellow!40!black}  
\rotatebox{20}{\RIGHTarrow}}}}} \hspace*{-2pt} #1
	
}
\end{noindent}
}

\newcommand{\sentenceexamples}[1]{
	
\vspace{1em}		
{\setstretch{1.0}	
#1
}
\vspace{.7em}		

\noindent\hspace{-4pt}}



\newcommand{\FeqMtimesA}{\hspace{-2pt}{\small \ensuremath{F}=\ensuremath{M}{\texttimes}\ensuremath{A}}}
\newcommand{\piCalculus}{\ensuremath{\pi}-calculus}


\newif\iffootnote
\let\Footnote\footnote
\renewcommand\footnote[1]{\begingroup\footnotetrue\Footnote{#1}\endgroup}

\newcommand{\AcronymText}[1]{{\iffootnote\begin{footnotesize}{\textsc{#1}}\end{footnotesize}%
\else\begin{OldFootnoteSize}{\textsc{#1}}\end{OldFootnoteSize}\fi}}


\newcommand{\NLP}{\AcronymText{NLP}}
\newcommand{\POS}{\AcronymText{POS}}

\newcommand{\Gardenfors}{G\"ardenfors}


\newcommand{\Haskell}{\AcronymText{Haskell}}
\newcommand{\C}{\AcronymText{C}}
\newcommand{\NP}{\AcronymText{NP}}
\newcommand{\Cpp}{\AcronymText{C++}}
\newcommand{\RDF}{\AcronymText{RDF}}
\newcommand{\Java}{\AcronymText{Java}}
\newcommand{\IT}{\AcronymText{IT}}
\newcommand{\IBMinc}{\AcronymText{IBM}}
\newcommand{\XDG}{\AcronymText{XDG}}
\newcommand{\AI}{\AcronymText{AI}}

\newcommand{\HPSG}{\AcronymText{HPSG}}
\newcommand{\CAG}{\AcronymText{CAG}}
\newcommand{\Lisp}{\AcronymText{Lisp}}

\newcommand{\TS}{\AcronymText{TS}}

\newcommand{\ThreeD}{\AcronymText{3D}}

%?\usepackage[inline]{enumitem}
\usepackage{enumitem}




\usepackage[font=small,labelfont=bf]{caption}

\usepackage{xcolor}

\definecolor{Bkg}{RGB}{250,245,252}
\newcommand{\leader}[2]{\hspace{#1}\colorbox{Bkg}{#2}}


\newcommand{\saying}[1]{\vspace{2ex}\noindent{%%
				\leader{2em}{\begin{minipage}{.38\textwidth}{\footnotesize #1}\end{minipage}}\vspace{2ex}}}

\newcommand{\sayingsrc}[1]{\vspace{0ex}\\\hspace{2pt} --- {\footnotesize #1}}

\renewcommand{\labelitemi}{$\blacklozenge$}

\newcommand{\itemmark}{\raisebox{-4pt}{\rotatebox{90}{{\Large $\bracevert$}}}}

\usepackage{tikz}
\usetikzlibrary{positioning}
\usetikzlibrary{shapes,snakes}

\newcommand{\visavis}{vis-\`a-vis}

%\sectionfont{\fontsize{12}{8}\selectfont}


\newcommand{\underlines}{{\normalsize \lowBlank\hspace{.1em}}}
%\newcommand{\underlines}{{\normalsize \lowBlank\lowBlank\lowBlank\hspace{.5em}}}

\newcommand{\XAcrossY}{\mbox{\ensuremath{x} \texttt{across} \ensuremath{y}}}

\newcommand{\Rel}{\raisebox{.5pt}{\texttt{{\OldFootnoteSize r}}}}

\newcommand{\XRelY}{\mbox{\ensuremath{x} \Rel{} \ensuremath{y}}}

\newcommand{\tinyurl}[1]{{\raisebox{2pt}{{\scriptsize \url{#1}}}}} 

\usepackage[colorlinks=true]{hyperref}

\colorlet{urlclr}{red!40!magenta!50!orange}

\hypersetup{
 urlcolor = urlclr,
 urlbordercolor = cyan!60!black,
 linkcolor = red!30!black,
 citecolor = orange!30!black,
 citebordercolor = yellow!30!black,
} 


\makeatletter
\Hy@AtBeginDocument{%
	\def\@pdfborder{0 0 0}% Overrides border definition set with colorlinks=true
	\def\@pdfborderstyle{/S/U/W .25}% Overrides border style set with colorlinks=true
	% Hyperlink border style will be underline of width 1pt
}
\makeatother


%\titlespacing{\section}{.25\textwidth}{*2.8}{*-.4}[5pc]


\newcommand{\p}[1]{
	
	\vspace{.575em}
	#1	
}

\renewenvironment{abstract}
{\small
	\begin{center}
		\bfseries \abstractname\vspace{-.5em}\vspace{0pt}
	\end{center}
	\list{}{%
		\setlength{\leftmargin}{.43in}% <---------- CHANGE HERE
		\setlength{\rightmargin}{\leftmargin}%
	}%
	\item\relax}
{\endlist}


\let\OldSection\section
\renewcommand\section[1]{
	\vspace{5pt}
	\OldSection{#1}
	\vspace{-4pt}
}

%?\let\OldSubsection\subsection

%?\renewcommand\subsection[1]{
%?	\vspace{-5pt}
%?	\OldSubsection{#1}
%?	\vspace{-16pt}
%?}

\usepackage{transparent}

\definecolor{logoCyan}{RGB}{66, 206, 244}
\definecolor{logoBlue}{RGB}{4, 2, 25}

\titleformat*{\subsection}{\Large\bfseries}

\let\OldSubsection\subsection
\renewcommand\subsection[1]{

	\vspace{12pt}
	
    %{\LARGE
	%\colorbox{logoCyan}{%
	%\begin{minipage}{\linewidth}
		\OldSubsection{% 	
			\hspace{-2.75em}
			%\begin{minipage}{}
			\protect\raisebox{-5pt}{%
			\colorbox{logoCyan!50}{\hspace{2.1em}}}%
			\hspace{-5pt}{\protect\transparent{0.3}{\colorbox{logoBlue!80}{\protect\transparent{1}{%
						   \protect\raisebox{1pt}{\textit{{\large #1}}} }}}}
			%\end{minipage}
		}
	%\end{minipage}
    %}
    %}
	\vspace{-6pt}
}


\newsavebox{\twolinebox}

\newcommand{\stwoline}[1]{%
\sbox{\twolinebox}{\raisebox{-3pt}%
{\parbox{7.4cm}{\linespread{1.25}\selectfont\raggedleft{\textit{\textbf{{\large #1}}}}}}}}

\newcommand{\twoline}{\usebox{\twolinebox}}

\newcommand{\subsectiontwoline}[1]{\stwoline{#1}\subsection{\twoline}}

\newcommand{\spsubsection}[1]{%
\vspace{-.25em}
\subsection{#1}
\vspace{.3em}
}
\newcommand{\spsubsectiontwoline}[1]{%
\subsectiontwoline{#1}
\vspace{.3em}
}




\let\OLDthebibliography\thebibliography
\renewcommand\thebibliography[1]{
\let\section\OldSection
\setlength{\leftmargin}{-4pt}
\vspace{.1em}
\OLDthebibliography{#1}
\vspace{.7em}
\OldFootnoteSize 
\setlength{\parskip}{0pt}
\setlength{\itemsep}{2.3pt plus 0ex}
\raggedright
}

\makeatletter
%\def\@biblabel#1{\hspace{-6pt}[#1]}
\def\@biblabel#1{\hspace{-6pt}#1}
\makeatother

\newcommand{\bibtitle}[1]{{\small \textit{#1}}}
\newcommand{\intitle}[1]{{\hspace{3pt}\textls*[-80]{\texttt{\textit{#1}}}}\hspace{-1pt}}

\renewcommand{\i}[1]{\textit{#1}}




\newcommand{\itemtitle}[1]{{\color{green!10!red!40!black} \textls*[-80]{\texttt{#1}}}}
\newcommand{\MIT}{\AcronymText{MIT}}

\newcommand{\nulltt}{\AcronymText{\texttt{null}}}
\newcommand{\Maybe}{\AcronymText{\texttt{Maybe}}}
\newcommand{\bind}{\AcronymText{\texttt{bind}}}
\newcommand{\return}{\AcronymText{\texttt{return}}}

\newcommand{\XML}{\AcronymText{XML}}
\newcommand{\XQuery}{\AcronymText{XQuery}}
\newcommand{\XQueries}{\AcronymText{XQueries}}
\newcommand{\API}{\AcronymText{API}}
\newcommand{\IR}{\AcronymText{IR}}
\newcommand{\Clang}{\AcronymText{Clang}}
\newcommand{\DAG}{\AcronymText{DAG}}
\newcommand{\DAGs}{\AcronymText{DAG}s}

\newcommand{\Qt}{\AcronymText{Qt}}
\newcommand{\NL}{\AcronymText{NL}}
\newcommand{\langC}{\AcronymText{C}}
\newcommand{\SDK}{\AcronymText{SDK}}

\newcommand{\SOV}{\AcronymText{SOV}}
\newcommand{\SVO}{\AcronymText{SVO}}

\newcommand{\YtwoK}{\AcronymText{Y2K}}

\newcommand{\qi}[1]{\q{\textit{#1}}}

\newcommand{\CSharp}{\AcronymText{C\#}}

\newcommand{\Tesniere}{Tesni\`ere}

\newcommand{\sentenceExample}[1]{``#1"}


\newcommand{\expose}{expos\'e}
\newcommand{\ExistsX}{\ensuremath{{\exists}X:Hugo(X){\wedge}Cat(X)%
{\wedge}Sleeps(X,on-the-couch)}}
\newcommand{\Exists}{\ensuremath{\exists}}
\newcommand{\ndim}{\ensuremath{n}-dimensional}

\usetikzlibrary{decorations.pathmorphing}

\newcommand{\longdash}{---}

\usepackage[breakable]{tcolorbox}% http://ctan.org/pkg/tcolorbox

\newenvironment{frquote}{%
%\begin{fcolorbox}{yellow!20!gray}{red!5}
%\begin{flushright}
\begin{tcolorbox}[
	colback=orange!6,
	colframe=yellow!20!gray,
	width=1.95\linewidth,
	arc=3mm, auto outer arc
	]
\begin{scriptsize}
\begin{minipage}{57em}
\begin{flushright}
\begin{minipage}{58em}}{%
\end{minipage}
\end{flushright}
\end{minipage}
\end{scriptsize}
\end{tcolorbox}
%\end{flushright}
}

\newcommand{\kErn}{\ensuremath{\mathcal{K}}}
\newcommand{\tSys}{\ensuremath{\mathbb{T}}}

%\newcommand{\tSys}{\ensuremath{T}}

\newcommand{\procOne}{\ensuremath{\mathcal{P}_1}}
\newcommand{\procTwo}{\ensuremath{\mathcal{P}_2}}
\newcommand{\fFunT}{\ensuremath{\mathcal{F}_t}}
\newcommand{\ChaOne}{\ensuremath{\mathcal{C}_1}}
\newcommand{\ChaTwo}{\ensuremath{\mathcal{C}_2}}
\newcommand{\CarOne}{\ensuremath{\mathfrak{c}_1}}
\newcommand{\CarTwo}{\ensuremath{\mathfrak{c}_2}}
\newcommand{\CarHandoffCar}{\ensuremath{\mathfrak{c}_1{\looparrowright}\mathfrak{c}_2}}
\newcommand{\ChpAlliedChx}{\ensuremath{\mathcal{C}_p{\circledcirc}\mathcal{C}_x}}
\newcommand{\Chp}{\ensuremath{\mathcal{C}_p}}
\newcommand{\Chx}{\ensuremath{\mathcal{C}_x}}
\newcommand{\ChaProductCha}{\ensuremath{\mathcal{C}_1{\otimes}\mathcal{C}_2}}

\newcommand{\stageLbl}{\ensuremath{\mathcal{L}}}
\newcommand{\DigammaR}{\ensuremath{{\varsigma}{\pointer}{\stageLbl}}}

\newcommand{\sigmaCalc}{\ensuremath{\varsigma} Calculus}
\newcommand{\voidP}{void*}
\newcommand{\ChaAppendCar}{\ensuremath{\mathcal{C}{\oplus}\mathfrak{c}}}
\newcommand{\LappendX}{\ensuremath{\mathcal{L}{\ll}\mathcal{X}}}
\newcommand{\Lprime}{\ensuremath{\mathcal{L}'}}
\newcommand{\LpreX}{\ensuremath{\mathcal{L}}}
\newcommand{\tTyp}{\ensuremath{\mathfrak{t}}}
\newcommand{\xy}{\ensuremath{xy}}



\begin{document}\title{Phenomenological Formalization: The Case of 
  Cognitive Grammar}
\p{Cognitive Grammar in relation to Linguistics 
is an interesting analog to Phenomenology in 
relation to Philosophy.  Thematically, Cognitive Grammar 
approaches linguistic phenomena from the perspective 
of language users' experience of percptual and 
situational structures \mdash{} as such, language 
structure is a kind of formalization, communication, 
and/or re-construction of patterns of organization 
within consciousness.  This perspective certainly seems 
amenable to Phenomenology more immediately than paradigms 
that connect language-meanings to logical 
propositions, or to the mechanical operations 
of semantic and grammatical rules.  Meanwhile, 
methodologically, Cognitive Grammar embraces a style of 
analysis grounded in first-person perspective; the linguist's 
reflective judgments on which potential expressions in a 
given language would be deemed acceptability to typical speakers 
of that language, and why.  True, that kind of subjective 
but evidence-based methodology \mdash{} linguists using their 
own sense of syntactic propriety and semantic coherence 
to determine qualifications like the acceptability of 
particular sentences, but doing so in the guise of generic 
speakers' language-use rather than any idiosyncratic 
preferences the linguist herself may have \mdash{} is common 
across many paradigms of linguistic research.  In the 
case of Cognitive Grammar, however, the researcher imagines 
linguistic phenomena situated in hypothetical perceptual and 
interpersonal contexts.  Issues like acceptability are 
accordingly assessed with consideration to the overall 
experiential, practical, and situational enacting and 
conceptualizing that would constitute the cognitive givens 
talked about and grounding communicative practice.  Insofar as 
practical and experiential context is treated as intrinsic amongv
the details needing to be consulted for trustworthy assesments 
of, say, acceptability \mdash{} rather than comparingb a given 
linguistic performance against dictionaries or style-books 
\mdash{} Cognitive Grammar eveals a metatheoretic commitmnt to the 
practical/experiential environing of cognition in general as 
something that can be systematically investigated, enough to make 
linguistics work as a rigorous science.
}
\p{Cognitive Grammar, in short, requires for the coherence of its 
scientific aspirations that a kind of first-personal but 
social-pragmatic analysis \mdash{} which we can argue is quite akin to 
Phenomenology \mdash{} can be performed which is also 
scientifically rigorous and amenable to third-person follow-up; 
phenomenological claims being reviwed and synthesized within a 
scholarly community.  Cognitive Grammar needs a story of 
its scientific merit that overlap with Phenomenology's 
parallel story: that first-person reports do not devolve into 
stream-of-consciousness but can mine experience for regulative 
patterns and recurring desiderata that can be critically 
examined by a sufficiently sympathetic research community, 
giving rise to theoris and models which can claim 
third-person objectivity, or at least claim that a reasonable 
disputative process exists to make claims toward objectivity 
evaluable and, as such, potentially believable.  In short, 
the phenomenologist can claim that an organizing gestalt she 
thinks is formulative in her experience as both consciously 
immediate and logicall unified, described from a reflective 
stance, is not a \q{subjective} phenomena in the epistemological 
sense of something idiosyncratic to her, but instead a 
first-personal phenomenon she can approach as something which 
others can find also in their own experience; similarly 
the cognitive linguist can consider her acts of producing or 
receiving linguistic artifacts, in representative 
perceptual-situational contexts, bearing 
in mind that she is a fluent speaker of one or more 
languages and can anticipate the language-processing practice 
of others.  The very fact of fluency, or competence, 
implies a certain degree of credibility in ascribing 
linguistic beliefs to others.  Accordingly, so long as 
we consider a speaker fluent is a language, we commit 
to finding her basically trustworthy in guessing how fellow 
speakers would respond to artifacts in that language, and 
therefore to generalize from subjective judgments outward.  
This is a good example of how furst-person reflection can be 
scientificaly operationalized for multi-party disputation; 
Phenomenology seems to envision a similar dialectic but 
one expanded outward from language to cognitive acts and 
experiential groundings in thir totality.   
}
\p{From at least a metascientific perspective, at least, a 
dialog between Phenomenology and Cognitive Grammar 
seems promising.  Both disciplines topicalize and fine-tune 
related but complementary methodological issues.  
Phenomenology, one the one hands, attends more vigorously 
to a philosophical problemmatics of origination and 
a \q{critique of metaphysics}, making a willful 
\q{bracketing} of extra-experiential assumptions an intrinsic 
part of its method.  Cognitive Grammar may be sympathetic to 
but not necessarily dependent on similar philosophical 
performances.  For example, an important issue in modern 
linguistics is whether we have \q{innate} mental faculties for 
creating and understanding language or whether language 
originates as a specialization or repurposing of some 
other, more general or primordial faculty.  Whatever our 
intuitions about that, surely those are the kind of 
assumptions that should be \q{bracketed} in a 
phenomenology-style analysis.  By comparison, while 
Cognitive Grammar appears to be internally consistent 
irrespective of any particular opinion \visavis{} 
linguistic innateness, I'm not aware of any crucial 
importance attached to either entertaining or suspending 
beliefs about innateness (or any other foundational 
linguistic-cum-philosophical subject matter, like whether 
a universal grammar underlies individual grammmars 
or what sort of preceptual inputs qualify as 
\i{linguistic} stimuli durin childhood).  Certainly 
Cognitive Grammar methodology is not \i{constituted} by 
the \i{absence} of such belief, the way that Phenomenology 
is essentially shaped by its counter-traditional, 
\q{Cartesian} style.  
}
\p{Notice however that Cognitive Grammar does intrinsically orient 
analysis to disciplinary commitments \mdash{} to the accuracy of 
subjective syntactic/semantics assessments and reports, to 
the value of hypthetical experiential episodes as proxies 
for actual ones, to the approximate commonality of experience 
among a language community \mdash{} which cannot be just passively 
accepted (or rejected); they must be actively (even if 
unconsciously) presupposed in the research process.  Phenomenology, 
given how it thematizes both the legitimacy of 
(experientially-grounded) presuppositions and the suspension 
of presuppositions drawn indirectly from scientific or 
philosophical beliefs, therefore points to a latent 
contrast in metalinguistic commitments between those 
that are tangential to and those that are constitutive 
of the Cognitive Grammar affiliated paradigms: 
without self-consciously witholding considerations 
on presumptive \q{metaphysical} linguistic matters 
(like innateness vs. linguistic empiricism), Cognitive 
Grammar does not necessarily to endeavor to classify 
it's practitioners beliefs on the axis of bracketing 
vs. disclosure \mdash{} the commitments that can be 
suspended partitioned from those that seem either 
experientially irrefutable or transcendentally prerequisite 
for any notion of experiential certainty in the first place.  
Scientists presumably feel that an \i{entirely} 
presuppositionless science is impossible, while a more 
philosophical angle can investigate how, nonetheless, 
presuppositions vary in their metaphysical posture.  
Belief in a \q{poverty of stimulus} suggesting innateness, say 
\mdash{} or, conversely, belief that the rich experiential structure 
surrounding even small-scale language acts presents substantial 
language-acquisitional stiumli (thereby falsifying the 
premise of poverty and that avenue of advocacy for 
innateness) \mdash{} can both be claimed to stand differently 
in the ebb and flow of cognitive-linguistic metatheory than 
commitments to, say, the disputational objectivization of 
first-person reports.   
}
\p{At the same time, Cognitive Grammar can be claimed to 
illuminate a wider, or at least a different, intellectual spectrum 
for the objectivizing dialectic than Phenomenology alone.  What 
I mean is that Phenomenology's path from subjective reflection 
to rationalistic claims is fairly repetitive: philosophers 
discuss specific hypothetical/prototypical episodes of consciousness 
and then write about them in an academic style which invites community 
reflection and to be read against certain scholarly traditions.  
Cognitive Grammar does all that too, but there are introduced by virtue 
of its linguistic subject-matter other discursive facets: the schematic 
representation of linguistic structure (e.g. parse trees); 
the diagramming of perceptual-situational gestalts; even technological 
applications...  Scolars can debate which tactics for 
representing language structure are most suited or most organically 
consistent for describing language artifacts via the 
Cognitive-Grammatic lense.  The representations specific to 
Cognitive Grammar can be contrasted with those developed 
in other grammar theories and their correlative 
research programs, such as Combinatory, Dependency, and 
and Head-Driven Phrase-Structural Grammars.  These disparate 
programs are, also, not exclusively focused on 
explanatory revelations about human language \mdash{} contemporary 
linguists also have practical concerns, like text mining 
and Natural Language Processing.  One question for 
semantic and syntactic theories is whether they are 
buttressed by computational evidence \mdash{} whether 
software created according to the language models 
and representations are successgul for useful tasks, 
like classifying documents, machine translation, 
or pulling information from Natural Language resources.  
}
\p{Cognitive Linguistics in general does not need to 
mimic these practical concerns \mdash{} after all, if language 
understanding does indeed build off of integrated 
percepual and experiential situatedness, we 
should be skeptical about the prospects for computers 
and \q{Artificially Intelligent} agents to parse language 
with anything like native fluency (even setting aside how  
a machine's language processing, correct or not in terms 
of our desired practical ends, cannot be equated to 
human empathic understanding).  Contrariwise, however, 
we can also speculate that if the gist of Cognitive Grammar is 
how a plethora of environing situations \mdash{} the 
spatial and force-dynamic arrangements trenchant to our construal 
of practical scenarios \mdash{} can be sorted into a 
relatively small set of canonical schema, corelated with 
details like noun tense and morphosyntactic modes of agreement, 
then a computational treatmnt of language need 
not percive the full spectrum of situational possibilities 
in full nuance, but only ascertain which one of several 
schema are evidenced in a given linguistic artifact.  
Cognitive Grammar would seem to support intuitions that 
truly humanlike \q{AI} Natural Language Processing is 
an unrealistic goal but that, at least potentially, computational 
NLP can be achieved to some useful approcimation, and employing 
cognitively realistic languag models can help that project.  
But I have no firm commitments in either direction; my point is 
that the technological and computational dimensions of 
linguistics are relevant for Cognitive Grammar whether 
or not this approach is sympathetic to the, at times, 
reductionistic paradigms of conventional 
formal/computational linguistics.  Unlike pre-communicative 
experience, language is structured by the rigors of 
coordinated thought and normative signification \mdash{} as 
language users we must shoehorn private 
experience into prototypes that can be expressed to others.  
Language, in effect, reveals that even rawly immediate 
experience has a logical order that, to some approximation, 
can be studied in abstraction \mdash{} first because 
language merely by existing shows that this abstracting 
does happen, in one fashion; and second because language is 
a platform for studying it.  Via language we can empirically 
study the articulations of conscious experience transcribed 
into a normative semiosis.
}
\p{Phenomenology, also, gathers consciousnss as reflected upon into 
something partially abstract, or with a structure that can be 
abstractly re-considered.  There are some parameters of organization 
that seem so primeval to experience as to be indubitable.  
One is the episodic nature of perception and of deliberate action; 
mental phenomena seem neither sliced into instantaneous moments nor 
streaming in expansive segments.  Instead, conscious reality is a chain 
of continuous but brief interludes: in one hypothetical kind 
of scenario, I open a door, which exercises one family 
of sensate and kinaesthtic attitudes; then I walk across th room; 
then I look out the window, and so forth.  Such episodes stick 
together in language and maybe in memory, if I have occasion 
to say, or to later recall, that \q{I walked into the room} \mdash{} the 
state of affairs I commit to memory or verbiage is actually a 
sequence of smaller episodes.  This in turn issubsume into 
larger units, like when a narration reads \q{I walked into 
my living room, sat at my desk, and began to read}.  What 
stands out inviolably from this episodic bricolage is 
\mdash{} minimally the temporality of, but more consequentially, 
the multiple scales of the temporality of consciousnss.  
Episodes link together and sequehtially coalesce into a 
more durational scale.  Phenomenology marks this scale-transition 
terminology; \i{protention} and \i{retention}, for instance, 
connoting the episodic unity of consciousness over short 
time spans, is contrasted with concpts like 
\i{memory} and \i{anticipation}.   
}
\p{Meanwhile, an equally salient structuring principle lies in 
the epistemic attitudes that are carried into perceptual/enactive 
episodes: consider the contrast between casually looking 
at the window expecting to see the same scene as 
any other afyernoon and looking inquisitively because I 
hear a sound outside.  Or maybe I see Hugo sleeping on the sofa, 
a \q{seeing} which is actually somewhat inferential, because he's 
always sleeping on the sofa this time of day.  
I don't need to confirm \mdash{} if I briefly glance and see him 
lying still and relaxed \mdash{} that the tabby cat on the couch is 
Hugo and that he's sleping, for me to think to myself (and maybe 
say to someone else) that \q{Hugo is sleeping} or \q{Hugo is 
sleeping on the sofa}.  Of course, such a scenario an 
be altered: if we have two tabby cats then I may want to look 
more discriminately to ascertain that it is Hugo on the sofa; 
or I may want to observe him for a few seconds to check 
\i{if} he is sleeping, which is a different perceptual 
event and a different conscious experience that passively 
seeing \q{Hugo sleeping} in the throes of an assuming 
that he is, in fact, sleeping, by analogy to how seeing 
daylight out the window is \i{passive} if it's mid-afternoon 
but may be a premeditatd investigation if I have just 
woken up and don't know the time.
}
\p{Here, too, is an apparently \q{transcendental} or \q{categorical} 
(allowing a Kantian echo) pattern in consciousness \mdash{} the sliding 
scale of activeness and passiveness in perceiving; the 
varying degree of epistemic thematization saturating 
different parts of experience.  There are things we take for granted 
based on the familiarity of our most typical environments, becoming 
part of the propositional hum of things generally disclosing 
themslves in ways that carry information but no surprises.  
There are other things we are consciously aware of wanting to know.  
Some of our information about the world comes from presumptive 
expectations that we only really register when they are 
violated \mdash{} say I suddenly see Hugo, having awaken, jump 
down from the sofa; now I have to update my baseline conception 
of how things are obtaining in my immediate surroundings.  
In that case I would not be startled to hear him scratching 
the chair, as compared to if I had not seen him wake up: 
our construals are woven from baseline assumptiona and then 
the occasional anomalies where my assumed picture must be revised.  
But layered within this cycle of passive experience and 
disruptive revision is another genre of epistemic episodes which 
is more deliberate and proactive; planning perceprual encounters 
around specific, epistemically thematized concerns, like 
looking out the window to trace an unfamiliar sound.  
This scale of epistemic passivity and activity can be 
set against \mdash{} and in consciousness interweaves with 
\mdash{} the scales of temporality.
}
\p{We have then at least this matrix \mdash{} the double-scale of temporal 
granularity and epistemic attentiveness \mdash{} which seems 
to emanate from the depths of consciousness and yet have something 
like an abstract structure.  Looking over the phenomenological 
(and philosophy of mind) tradition, we can find other such 
prereflective abstracta, like the contrast between hyletic sense-data 
(the purely embodied qualitative character of red or of silky 
smoothness) and qualia as propositional content 
(\i{red} and \i{smooth} as predicates); or the contrast between 
fully private experience (like corporeal sensations) and 
intersubjectively negotiated perception \mdash{} we may not see the same 
scenes in exactly the same ways, but at least in productively 
\i{similar} ways, if we are in similar locales and vantage-points.  
A full catalog of these structures would involve 
a historical overview of Phnomenology as a whole, which is 
divergent from my aims in this paper.  I want to highlight, 
however, that Phenomenology produces a philosophical narrative about 
certain ur-structures \mdash{} described, if all goes well, not 
from prior scientific or world-view commitments but from 
commonsensical meditations on (mundane, everyday) consciousness   
\mdash{} which can be thematized with the extra performative 
precision of scholarly discourse; with the care of finding 
the right language, entrenching usages, diagramming the 
conceptual relations which should be easier 
to memoize precisely because of terminological rigor 
(\q{protention}; \q{episodic}; \q{epistemic 
attitude}; \q{enactment}).  These structurs get lifted out 
of the experiential realm and begin to condensate into 
an abstract regime, something defended like intellectual 
turf, discussed and analyzed.
}
\p{Cognitive Grammar also has ur-structures \mdash{} for sake of 
argument, consider on one hand the linguistic 
double-parameters of spatial and force-dynamic schema and, 
on another, the phenomenological double-scales of temporal 
granularity and epistemic thematization.  
Once entered in the intellectual ledger, the cognitive-linguistic 
schema become subject to a range of formal systemizations, 
from technical representations of language structures 
to lexicons, compilations of parts of speech and type-theoretic 
lexical classifications as well as of inter-word 
relations, even computer code for representing 
or automatically building representations of language artifacts.  
There have not been thus far analogous options for 
elaborating the phenomenologist's ur-structures 
\mdash{} no code-repositories curated by Phenomenology 
research groups, no technological enhancements 
of phenomenological literature, for the most part.  
This is doubtless a manifestation of the historical 
context where Phenomenology has emerged, but we are 
now in a different context.  Perhaps some 
21st-century phenomenologists will engage with computers 
as agressively as Husserl engaged with mathematics in 
the \i{Formal and Transcendental Logic} or the 
\i{Investigations}.
}
\p{Language, in short, is \mdash{} at least on one level \mdash{} a 
\i{formal system}, notwithstanding that it 
is also a human medium possessing an arguably irreducible 
layer of human subtlety and contxt-sensitivity.  Linguistics 
in general, and perhaps Cognitive Grammar in particular, 
has to bridge experience and formalization.  Spatial and 
force-dynamic schema are \mdash{} for example and so to speak \mdash{} 
one one side of the bridge as perceptual gestalts; but their 
formal trace is excavated analytically through syntactic 
and morphosyntactic theories; and a systematic model 
of morphosyntax involves something like formal 
language-representation.  Insofar as parsing structures are 
described via \q{parse graphs} \mdash{} the parsed sentence notated, 
with the supplemental data marking the contrast between a 
sentence spoken in an unstudied human context and the sentence as an 
object of analysis, as a graph of inter-word relations 
\mdash{} then morphosyntactic agreement is a paramter manifest 
in specific inter-word pairs.  So morphosyntactic agreement 
qua linguistic phenomenon straddles the experiential 
and formal realms \mdash{} it can be studied in terms of 
what perceptual situations call for a particular 
noun or verb to be aligned with a particular verb, adjectuve, 
or adverb; it can also be studied as an enrichment of a 
formal structure.  Through such formal strata 
Cognitive Grammar is not siloed in its study of perceptual 
groundings, but widens its disciplinary circle to reach a 
constellation of formal techniques that resonate 
with grounding experiential intuitions in suggestive ways, 
revealing both the potential and limits of formalization.  
}
\p{I would like to juxtapose this interdisciplinary continuation 
of Cognitive Grammar with the academic dissemination of 
Phenomenology \mdash{} how the professional philosophical milieu where 
Phenomenology is mostly practiced shapes its 
transition from first-person speculation to an intellectual 
enterprise which, at least in its ways of organizing social 
resources, operates as a science.  Cognitive Grammar operates 
in a similar milieu, but the strata of formalization 
evident in linguistics adds a wrinkle that has no 
direct analog in philosophy.  Alongside the discursive and 
institutional norms which guide the systematization 
of Cognitive Grammar \mdash{} the discipline of 
acdemic writing, peer review, conferences \mdash{} there are 
also structural norms sited in linguistic formalizations 
through which cognitive theories can be explored.  
This is a development which phenomenologists should observe 
carefully, because the structural elaboration upon 
phenomenological research can be augmented beyond 
\i{just} inter-textual disputation in an academically curated 
speech situstion, and can engage with formal 
systems as logical and technological artifacts.  
Scholars, of course, debate, refine, and mathematize 
formal systems in similarly self-conscious academic 
circles; but insofar as they are modeled and 
simulated in technological and computational environments, 
formal systems \i{also} become technical artifacts 
which can be explored and manipulated.  Alongside the 
\q{Communicative Rationality} of academic performance, they 
buttress a \i{technical} rationality of 
implementing computer models, crafting the algebra of 
formal systems, codifying inter-system translations, 
and so forth.  Computational linguistics, 
to take one example, is a technical as well as an 
intellectual practice.  
}
\p{The paradigms of these research programs that emphasize 
engineering over academic disputation don't necessarily 
align seamlessly with Cognitive Grammar, and still less 
with Phenomenology.  Despite their \i{internal} commitment 
to \q{bracketing} phenomenologists are still educated 
participants in the modern world and many presumably do 
accept in broad outlines the scientist's worldview.  
Many phenomenologists probably \q{unofficially} 
believe that mental phenomena have mundane neurophysical 
explanations (they are not magic or 
divine revelations or flux in a sacred ether), 
however opaque to consciousness itself.  So Phenomenology 
is not \i{constitutively} antagonistic to a \q{natural science} 
of cognitiona and perception.  In the Philosophy 
of Mind, moreover, the path to materialistic perspectives on 
mental phenomena seems to diverge in two irections \mdash{} one being 
the direct study of material systems themselves that we reasonably 
believe are seats for the neurophysical correlates of consciousness, 
like neurons and synapses; the other is a mkre functioanlist 
attemp to describethe systematic organization of 
such material systems, on the premise that it is easier to make 
scientific process by seeing the brain (or the entire 
embodied nervous system) as a large-scale functional system than 
by reductively studying the biology and physics of 
microscale constituents, like nerce cells.  Sure, we may believe that 
vision is driven by cells in the eye and optic nerve, as 
well as specific brain regions; but the explanatory gap between 
whatever organic properties we may discover researching these 
bodies and the lived immediacy of visual qualia seems no 
less expansive.  At least as a supplement to such microphysical 
investigation, functionalist methodologies can potentially 
narro, even if not elimibate, such explanatory gap.  In Cognitive 
Grammar, for instance, the subtle variation between schematically 
similar situations is in a sense transcendended by morphosyntactic 
rules; our language-forming process of mapping perceptual 
givens to morphosyntactic and lexical prototypes and 
finding a communicatively stable encoding for them 
certainly seems amenable to functional description: such 
meta-cognitive finessing of our preceptual surround is a 
good candidate of a \i{functionality} available to our minds,  
as we are (and in our being) intelligently intersubjective 
and information-processing life-forms.  Insofar as 
we have formal presentations of perceptual and enactive 
structures that seem both subjectively realistic \i{and} 
faithful to a kind of inner logic, we have the possibility for 
a rapproachment between Phenomenology and cognitiv 
functionalism, the two paradigms bridged by 
the \i{functional} utility of cycling between 
immediate and subtly particular experience and 
formally tractable concept-systems we us to construct 
productive precis of our environing situations 
\mdash{} which in turn can be modeled as formal systems 
and investigated in that light.  Language bears witness 
to this experiential-to-formal-to-experiential rotation 
in an especially ubiquitous and well-structrued topos, 
which is why linguistic methodologies like Cognitive Grammar 
can be a useful case-study for phenomenological formalization 
in general. 
}
\p{The preceding discussion has set forth why I believe the 
formalizing strategies available to Cognitive Grammar should be 
interesting to Phenomnology, in part because they are perforce 
available to Phenomnology as well.  A formalized cognitive 
linguistics may not be immediately a formalized Phenomenology, 
but but it comes tantalizingly close to that, as I hope 
to show before the end of this paper.  But I have until now 
been addressing this formalization obliquely, talking 
\i{about} Phenomenology and Cognitive Grammar as scholarly 
phnomena rather than \i{within} either (or both).  Here, 
for most of what follows, I will focus in on 
one specific formal development of Cognitive Grammar 
from both a cognitive and computational perspective, 
which I will then conclude by placing in a more 
classically phenomenological context.
}
\p{}
\section{Varieties of Hypergraphs}
\p{Any notion of hypergraphs contrasts with an 
underlying graph model, such that some element treated as 
singular or unstructured in regular graphs becomes a multiplicity 
or compound structure in the hypergraph.  So for example, an 
edge generalizes to a hyperedge with more than two incident 
nodes (which means, for directed graphs, more than 
one source and/or target node).  Likewise, nodes might generalize 
to complex structures containing other nodes (including, 
potentially, nested graphs).  The \q{elements} of a graph 
are nodes and edges, but also (for labeled/weighted 
and/or directed graphs) things like labels, weights 
(such as probability metrics associated with edges), and directions 
(in the distinction between incoming and outgoing edges incident to any node).  
Potentially, any of these elements can be transformed 
from a unit to a plural structure, a process 
I will call \i{diversification}.  That is, 
a hypergraph emerges from a graph by \q{diversifying} 
some elements, rendering as multiplicities what 
had previously been a single entity \mdash{} nodes become node-sets, edges 
become grouped into larger aggregates, labels generalize to 
complex structures (which I will call \q{annotations}), 
etc.  Different avenues of diversification give 
rise to different varieties of hypergraphs, as 
I will review in the next several paragraphs.
}
\p{\begin{description}\item[Hyperedges]  Arguably the most common model of 
hypergraphs involves generalizing edges to hyperedges, 
each of which (potentially) connect more than two nodes.  
For directed graphs,  
directed hyperedges have a \q{source} node-set and a \q{target} node-set, 
either of which can (potentially) have more than one node.  
Note that this is actually a form of node-diversification 
\mdash{} there is still (in this genre of hypergraph) just one 
edge at a time, but its incident node set (or source and 
target node-sets) are sets and not single nodes.  Another way 
of looking at directed hyperedges is to see source and target 
node-sets as integral complex parts, or \q{hypernodes}.  
So a (directed) hypergraph with hyperedges can also 
be seen as generalizing ordinary graphs by replacing 
nodes with hypernodes (such that one hypernode 
spans a \i{set} of \q{normal} nodes).
 
\item[Recursive Graphs]  Whereas hyperedges embody a relatively 
simple node-diversification \mdash{} nodes replaced by node-sets 
\mdash{} so-called \q{recursive} graphs allow compound nodes (hypernodes) 
to contain entire nested graphs.  Edges in this case can still 
connect two hypernodes as in ordinary graphs, but the 
hypernodes internally contain other graphs, 
with their own (on another \q{level}) nodes and edges.
 
\item[Link Grammar]   Since \i{labeled} 
graphs are an important model for computational models, 
we can also consider generalizations of labels to be 
compound structures rather than numeric or string labels 
(or, as in the Semantic Web, \q{predicate} terms,  
drawn from an Ontology, in \q{Subject-Predicate-Object} 
triples).  Compound labels (which I will generically call 
double- or multiple-annotations) are encountered 
in several different branches of mathematics and other fields.  
\vspace*{.3em}\\\hspace*{3em}
In particular, compound annotations can represent the 
rules, justifications, or \q{compatibility} which allows 
two nodes to be connected.   In this case the annotation 
may contain information about both incident nodes.  
This general phenomenon (I will mention further examples 
below) can be called a \q{diversification} 
on \i{labels}, transforming from labels as single units 
to labels as multi-part records.
\item[Hypergraph Categories]  Directed hyperedges have source- 
and target-\i{sets} of nodes (possibly ordered).  Allowing 
these sets to be \i{empty} yields structures similar to 
hypergraph categories in recent mathematical treatments, such 
as \cite{BrendanFong}, \cite{AleksKissinger}.  Hypergraphs 
in this characterization might be used to study software applications 
and real-time monitoring systems.  Hyperedges would then 
mark information flows between different sites \i{inside} the 
system.  A directed hyperedge \i{without} a source node-sets then 
models information \i{entering} the system \i{ab initio}.  
Likewise hyperedges without target node-sets could model 
empirical effects produced by the system, which from one 
perspective embodies information \q{leaving} the system.
\vspace*{.3em}\\\hspace*{3em}
Analogous models can be realized by, in lieu of 
hyperedges \i{without} source or target node-sets at all, 
designating special hypernodes representing empty 
node-sets; or hypernodes providing greater detail about 
origin- or destination-points for paths modeling 
information flows.  In software and some CyberPhysical 
systems, information \q{enters} the system via 
\q{events} \mdash{} that is, \i{events} are introduced as a 
primitive concept alongside \i{procedures}, though 
events themselves are not \q{implemented} (they have no 
\q{body}).  While procedures call other procedures, 
events \i{trigger} calls via event-to-procedure 
\i{connections}.   
Events which are not generated by components \i{in} the 
system then correspond to hyperedges without source-nodes, 
but the situation can be more precisely modeled via a 
special class of event hypernodes (analogous to 
events, as a distinct signature-bearing element from 
procedures, in event-driven programming; 
cf. \cite{JenniferPaykin}, 
\cite{PaykinKrishnaswami}, \cite{WolfgangJeltsch}). 
\item[Edges Incident to Edges]  A further generalization 
allows edges to \q{point at} other edges 
\cite[p. 10]{BalintMolnar}; \cite[p. 13]{BenGoertzel}.
The rationale for concordant constructions is that 
an edge \mdash{} qua relation between two or more 
components \mdash{} is itself a fact, datapoint, or 
assertion, and as such can be on some perspectives a 
singular element in a space of information.  
The Semantic Web, for example, includes \q{reified} 
edges (Subject/Predicate/Object triples).  
Reification permits properties to be stated about edges, 
such as provenance data (who asserts the edge and on 
what evidence) and context (the assertion is believed true 
at a specific time or in a given context).  
It is plausible then to allow reified edges 
to be treated as nodes incident to other edges.  As 
with hyperedges possessing empty source or target 
sets, the associated paradigm can be interpreted through 
a special class of hypernode: we can certainly 
consider nodes which refer to edges (perhaps via the 
latter's label or annotation).  Such \q{referring} 
is not necessarily an explicit further edge in graphs 
but could instea be \q{semantic} \mdash{} that is, an edge being the 
\q{value} associated with a node (or likewise  
one value in a multi-element hypernode).
\item[Channels]   Hyperedges, which connect multiple 
nodes, are still generally seen as single edges 
(in contrast to multigraphs which allow multiple 
edges between two nodes).  Analogous to the grouping 
of nodes into hypernodes, we can also consider structures 
where several different edges are unified into a larger 
totality, which I will call a \i{channel}.  
In the canonical case, a directed graph can group 
edges into composites (according to more fine-grained criteria 
than just distinguishing incoming and outgoing edges) which share a 
source or target node.  The set of incoming and/or outgoing  edges 
to each node may be partitioned into distinct 
\q{channels}, so that at one scale of consideration the network 
structure can be analyzed via channels rather than via 
single edges (this shows why channels mark edge \i{diversification} 
as well as edge \i{aggregation}, because the conceptual role 
of edges is potentially replaced by channels, which are composed 
of numerous edges).
\vspace*{.3em}\\\hspace*{3em}
As a concrete case, 
consider graphs used to model Object-Oriented programming 
languages, where a single node can represent a 
single function call.  The incoming edges are then 
\q{input parameters}, and outgoing edges are procedural 
results or \q{outputs}.  In the Object-Oriented paradigm, however, 
input parameters are organized into two groups: in addition 
to any number of \q{conventional} arguments, there is 
a single \this{} or \self{} object which has a distinct  semantic 
status (\visavis{} name resolution, function visibility, and 
polymorphism).  This calls, under the graph model, for splitting 
incoming edges into two \q{channels}, one representing regular 
parameters and a separate channel for the distinguished or 
\q{receiver} object-value.
\end{description}
}
\p{In each of these formulations, what makes a graph \q{hyper} is the presence of 
supplemental information \q{attached} to parts of an underlying (potentially 
labeled, directed) 
graph.  For a concrete example, consider a case from linguistics \mdash{} specifically, 
morphosyntactic agreement 
between grammatically linked words, which involves details 
matched between both \q{ends} of a word-to-word \q{link} 
(part-of-speech, plural/singular, gender, case/declension, 
etc.).  According to the theory of Link Grammar, 
words are associated with \q{connectors} (a related useful terminology, 
derived more in a Cognitive-Linguistic context, holds that 
words carry \q{expectations} which must be matched by other words they 
could connect to).\footnote{Link Grammar proper was developed by Davy Temperley and Daniel Sleator 
\cite{TemperleySleator}, \cite{GoertzelPLN}, and computer 
scientists such as Kenneth Holmqvist and Matt Selway have adopted 
concepts from Link Grammar as part of projects formalizing 
Cognitive Linguistics in the tradition associated (notably) with 
Ronald Langacker \cite{HolmqvistDiss}, \cite{MattSelway}  
The OpenCog system is an example of Link Grammar integrated with 
hypergraphs concerned with data models and procedural 
execution, embodied in the AtomSpace data management component 
\cite{GoertzelEtAl}, \cite{RuitingLian}, \cite[pp. 55ff]{GoertzelP2}.  
}  The word \i{many}, for instance, as in \i{many dogs}, 
carries an implicit expectation to be paired with a plural noun.  The actual 
syntactic connection \mdash{} as would be embodied by a graph-edge when a graph 
formation is employed to model parse structures \mdash{} therefore depends on 
both the expectations on one word in a pair (whichever acts as a modifier, 
like \i{many}) and the \q{lexicomorphic} details of its \q{partner} 
(\q{dog}, as a lexical item, being a noun, and \i{dogs} being 
in plural form).
}
\p{In Link Grammar terminology, both the expectations on 
one word and the lexical and morphological state of a second are called 
\q{connectors}; a proper linkage between two words is then a \i{connection}.  
For each connection there is accordingly two sets of relevant information, 
which might be regarded as a generalization on edge-labeling wherein edges 
could have two or more labels.  Furthermore, the assertion of multi-part annotations 
on edges permits edges themselves to be grouped and categorized: aside from several 
dozen recognized link varieties between words (which can be treated as conventional 
edge-labels drawn from a taxonomy, consistent with ordinary labeled graphs), 
edge-annotations in this framework mark patterns of semantic and syntactic agreement 
in force between word pairs (not just foundational grammatic matching, like 
gender and number \mdash{} singular/plural \mdash{} but more nuanced compatibility at the 
boundary between syntax and semantics, such as the stipulation that a noun in a 
locative position must have some semantic interpretation as a place or 
destination).  Insofar as these agreement-patterns carry over to other word-pairs, 
annotations mark linguistic criteria that tie together multiple edges 
in the guise of signals that a specific parse-graph (out of the space 
of possible graphs that could be formed from a sentence's word set) is correct.
}
\p{Ordinary directed graphs 
(not necessarily hypergraphs) already have some sense of grouping edges together, 
given that incoming and outgoing edges are distinguished; 
but this indirect association between edges does not 
internally yield a concordant grouping of the nodes 
at the sources of edges all pointing to one target node 
(or analogously the targets of one source node).  
In the theory of 
hypergraph \i{categories}, hypernodes come into play insofar as representation 
calls for the nodes \q{across} (i.e., at the other end of) 
incoming/outgoing edges to be pulled together.   
We distinguish an \i{incoming node-set} from an \i{outgoing node-set}; 
then we might treat these sets as integral \q{bodies} of inputs 
(respectively, outputs):   
\begin{cquote}The term hypergraph category was introduced recently ... in reference
to the fact that these special commutative Frobenius monoids provide precisely the
structure required for their string diagrams to be directed graphs with \sq{hyperedges}:
edges connecting any number of inputs to any number of outputs. ... 
We then think of morphisms in a hypergraph category as hyperedges 
\cite[p. 13]{BrendanFongThesis}.
\end{cquote}  
For a general mapping of  
categories to graphs where edges represent morphisms, this represents a 
generalization on the notion of \i{edges} themselves.  Suppose morphisms are 
intended to model computational procedures in a general sense (say, as 
morphisms in an ambient program state).  Because procedures can have any 
number (even zero) of both inputs and outputs, this implies a generalization 
wherein directed edges can have zero, one, or multiple source (and respectively 
target) nodes.  A corresponding hypergraph form is one where hyperedges have 
source- and target- hypernodes, but each hypernode models variant-sized sets 
of further \q{inner} nodes (or \q{hyponodes}); here the empty set can be 
a hypernode with no hyponodes.  Elsewhere, apparently an equivalent 
structure is called a \q{trivial system}:
\begin{cquote}Monoidal categories admit an elegant and powerful graphical notation 
[wherein] an object $A$ is denoted by a wire [and] [a] 
morphism \fAB{} is represented by a box.  The trivial system 
$I$ is the empty diagram.  Morphisms \uIA{} and \vAI{} 
... are referred to as \textbf{states} and \textbf{effects} 
\cite[p. 4]{InteractingConceptualSpaces}.
\end{cquote}
The system $I$ \mdash{} which we can also see as a hypernode with an empty 
(hypo-)node set \mdash{} may embody a procedure which has no internal 
algorithmic or calculational structure (at least relative to the 
domain of analysis where we might represent computer code).  In 
CyberPhysical Systems, a function which just produces a 
value (with no input and no intermediate computations) can also 
be called an \i{observation}, perhaps a direct reading from a 
physical \i{sensor} (accordingly, as in the above excerpt, a \i{state}).  
Dually, a procedure which performs no evident calculation and produces 
no output value, but has a CyberPhysical \i{effect}, can be called an 
\q{actuation}, potentially connected to a CyberPhysical \i{actuator} 
(an example of a sensor would be a thermostat, and an example of an 
actuator would be a device which can activate/deactivate a furnace 
and/or cooling system).
}
\p{In these examples I have cited hypergraphs in a linguistic (Link Grammar) 
and a mathematical (hypergraph categories) context.  These 
share some parallels insofar as a core motivation is to generalize and add 
structure to edges, either freeing edges from limts on 
node-arity (even allowing edges to be \q{unattached}), 
or supplying edges with structured (potentially multi-part) annotations.
}
\p{A somewhat different conception of hypergraphs is found in database 
systems such as HypergraphDB.  Graphs in that context express what have 
been called \i{recursive} graphs, wherein a hypernode \q{contains} or 
\q{designates} its own graph.  The basic idea is that for each (hyper-)node we 
can associate a separate (sub-)graph.  This can actually work two 
ways, yielding a distinction between (I'll say) \i{nested} 
and \i{cross-referencing} graphs.  In a \i{cross-referencing} graph, 
subgraphs or other collections of graph elements (nodes, edges, and/or 
annotations) can be given unique identifiers or designations and, as 
a data point, associated with a separate node.  Consider a case where 
nodes refer to typed values from a general-purpose type system; insofar 
as subgraphs themselves may be represented as typed values, a node could 
reference a subgraph by analogy to any other value (textual, numeric, 
nominal/enumerative data, etc.).  Here the (hyper)node does not 
\q{contain} but \i{references} a subgraph; the added structure involves 
subgraphs themselves being incorporated into the universe of values 
which nodes may quantify over.  Conversely, \i{nested} hypergraphs 
model hypernodes which have other graphs \q{inside} them, thereby creating 
an ordering among nodes (we can talk of nodes at one level belonging 
to graphs which are contained in nodes at a higher level).  Such constructions 
may or may not allow edges across nodes at different levels. 
}
\p{Note that nested hypergraphs can be seen as a special case of cross-referencing 
graphs, where each hypernode \nodeN{} is given an index \ix{}, 
with the restriction that when  \nodeN{} designates (e.g., via 
its corresponding typed value) a subgraph \sg{}, all of 
\sg{}'s hypernodes have index \ixminusone{}.  Cross-referencing 
graphs, in turn, can be seen as special cases of an overall 
space of hypergraphs wherein hypernodes are paired with 
typed values from some suitably expressive type system \tys{}.  
If \tys{} includes higher-order types \mdash{} especially, lists and 
other \q{collections} types which become concrete types in 
conjunction with another type (as in \i{list} qua generic 
becomes the concrete type \i{list of integers}), then 
hypernodes acquire aggregate structure in part by 
acquiring values whose types encompass multiple 
other values.  I will use the term \i{procedural} hypergraphs 
to discuss structures that model node diversification 
via mapping hypernodes to collections-types 
(with the possibility for hypernodes 
to \q{expand} or \q{contract} as values are inserted into or 
removed from the collection).
}
\p{Cross-referencing also potentially introduces a variant derivation of 
hypergraphs which proceeds by accumulating graph elements (nodes, 
edges, and labels or annotations) into higher-scale posits, rather 
than defining inner structures on elements.  Specifically, as a 
complimentary operation to diversification, consider \q{aggregation} 
of graph elements: the option to take a set of (hyper)nodes, 
(hyper)edges, and/or annotations as a typed value which can then 
be assigned to a (hyper)node.  In such a manner, higher-level 
structures can be notated with respect to graphs, which is one 
way to model phenomena such as \i{contexts} \mdash{} the kind of multi-scale 
patterns that are considered endemic to practical domains 
like the Semantic Web.  While it is informally acknowledged that 
a single-level interpretation of the Semantic Web is misleading 
\mdash{} the Semantic Web is not an undifferentiated mesh of connections, but 
rather an aggregation of data from many sources, which implies the existence of 
localization, contextualization, and other \q{emergent} structure \mdash{} 
there is no definitive protocol for actually representing this emergent 
structure.  This issue, in turn, is one of arguments 
for hypergraphs in lieu of ordinary graphs as general-purpose 
data representations. 
}
\p{The multi-scale, contextualized nature of the Semantic Web also points toward 
a conceptual duality in \i{how} graphs represent data.  On the one hand, graph 
structures \mdash{} especially in the case of the Semantic Web, which builds off of 
internet technology in general \mdash{} represent \i{relationships} between points 
or structures of data in some sense; in familiar web terms, the relata linked 
by graph edges are often \q{resources}, designated by unique web addresses.   
A more theoretical model might take the information \q{residing} at 
graph nodes as typed values.  But in either case a given node may stand in 
for an aggregate of information \mdash{} a single web resources may contain a 
theoretically unlimited supply of data, and a typed value can be of a list 
or tuple type internally containing its own body of information.  
Consequently, the full stock of data embodied in a graph may not lie 
primarily in the graph structure itself, but rather distributed among 
its nodes (that is, among nodes' associated data). 
}  
\p{Conversely, graph structures (with suitable semantic specifications) are  
also considered to be media for serializing arbitrary data structures, which 
implies representing all details, at all levels of hierarchical organization, 
via graph structures.  Insofar as nodes embody their own information spaces, 
such internal structuration must then be mapped to their own graphs, 
as part of a workflow to project arbitrary structured data onto a canonical 
format.  The theoretical corollary to this idea would be that node data 
(via associated typed values) has its own internal representation; 
that it is, every typed value has a corresponding graph structure 
that may be \q{contained} within a higher-scale node.  
}
\p{Different varieties of hypergraph forms complicate this picture 
because the structuring elements of hypergraphs include 
aggregate data within nodes as well as the space of 
edges and incidence relations.  Given the structures I have referred to 
as \q{procedural} hypergraphs, nodes can encompass multiple
internal values so long as they have a suitably well-defined 
internal structure.  We can analyze these possibilities 
by defining, for each hypernode, an \q{interface} or list of 
operations available for updating hypernodes' associated values.  
Which operations are proper depends on a hypernode's type: 
hypernodes associated with a single unstructured value 
should have one basic update operation, while nodes with 
list-like types would have operations to insert (and remove) 
values at different positions.  In general, procedural 
hypergraphs should differentiate between \i{atomic} hypernodes 
with one associated value; \i{tuple} hypernodes with a 
fixed array of values (potentially of multiple types); 
and \i{collections} hypernodes encompassing lists of 
values subject to append/insert/remove actions 
(in which case inserted values' types should 
have a predictable pattern; 
consider key- and value-type alternation in 
associative arrays).
} 
   
\p{Separate and apart from operations modifying hypernodes' values, there 
are also conventional graph operations \mdash{} adding and removing 
nodes and edges.  In combination, the graph-oriented and node-oriented 
operations present a variegated interface for manipulating 
procedural hypergraphs.  Such an interface then serves as a rigorous 
characterization of the overall hypergraph model \mdash{} the structuring 
elements expressed by enumerating graphs' transformation operations 
represent the particular features of each specific hypergraph 
variety.  In the case of procedural hypergraphs, many of these transformations 
are not graph-related per se but derive from hypernodes' collections or tuple 
types.  I contend that this is a useful property of procedural hypergraphs 
for reasons I alluded to in the introduction \mdash{} hypergraphs (or at least the 
data represented with them) need to work in a variety of computational contexts.  
The relatively unstructured form of graph data is not always appropriate 
from one context to another; the list-of-values or value-tuple structures 
embodied as hypernode data may be more consistent with internal representations 
in database or language-runtime engines, for example.  Procedural hypergraphs 
are appropriately flexible in that some data is modeled at the graph level 
proper while other data is modeled as lists, tuples, and similar data structures 
within individual nodes.  
}
\p{In many practical contexts graph structures are not implicitly used at all; 
the importance of Semantic Web-style representation is for intermediary 
structures, where information is routed among different environments 
(database, applications, serialization, and so forth).  Transformations 
between hypergraphs can be a central process in generic transformations 
between data structures proper to different contexts.  In effect, where 
there is a general need for data tranforms in routing between 
contexts \mdash{} e.g., database to application runtime \mdash{} we can hope to 
capture most of the relevant transform logic as specifically 
mappings between hypergraphs, with each hypergraph possessing a structure 
optimized for being initialized from one context (or for generating data 
used in another context).  This progression may involve restructuring 
wherein information modeled at the \q{inter-hypernode} level \mdash{} 
which we can call \i{hypernode} data \mdash{} tends to be migrated 
to the level \i{inside} hypernodes, which we can call \i{hyponode} data.  
In broad outline, hypergraph transforms can progress from relatively 
casual structures (with a preponderance of information expressed as 
\i{hypernode} data) to more constrained structures \mdash{} mapping 
hypernode to hyponode data, i.e., mapping data from structures 
\i{between} hypernodes to those \i{within} hypernodes \mdash{} where 
the hyponode data is regulated by hypernodes' types.
}
\p{In this section I identified different additions through which graph 
structures generalize to hypergraphs; a general-purpose hypergraph 
engine would need to support each of these variations, which entails 
enabling the complete repository of transform-operations applicable to 
different hypergraph varieties.  This includes generalizing 
edges to hyperedges by \q{diversifying} nodes to encompass mutiple values; 
representing nodes' internal structure in terms of data structures 
such as lists, tuples, and nested graphs; and generalizing edge-labels 
to annotations which may have multiple parts.  I have not yet discussed 
the possibility of grouping hyperedges into higher-level structures 
(what I called \q{channels}), but I will return to 
those details in a later section.  
}
\p{For the remainder of this paper, I will attend especially to 
hypergraphs modeling (and subsequently executing) computer code, 
because constructing a working runtime engine which runs source 
code, via hypergraph intermediaries, demonstrates a variety of concepts 
applicable to hypergraphs in general.  I will discuss a workflow 
connecting parsers, a runtime \q{virtual machine}, intermediate representations 
for hypergraphs, and inter-graph transforms.  Collectively 
this workflow implicates many of the capabilities which would be 
requisite for a general-purpose hypergraph software ecosystem.  
}
\p{Interested readers who would like to observe a concrete 
unfolding of this workflow are invited to download the 
code base accompanying this paper, where readers can examine 
the operations of parsers and code generators working with 
a built-in, general-purpose hypergraph library.  The demo  
is fully integrated with Qt Creator, a C++ Integrated Development 
Environment, and has no further dependencies (assuming users 
have a working Qt and C++ compiler; Qt is a popular C++ application-development 
framework).  The dataset includes instructions for 
experimenting with the hypergraph library via Qt Creator and examining 
runtime structures by executing demonstration scripts in Debug mode. 
}

\section{Truth-Theoretic Semantics and Enaction}
\p{Corrolary to the idea that roles often determine 
concepts, is the recognition that 
we tend to logically evaluate situations in 
functional terms, through the lens of what 
we (or any of our peers) are \i{doing}.  
Suppose my friend says this, before and after: 
\begin{sentenceList}\sentenceItem{} Can you put some almond milk in my coffee?
\sentenceItem{} Is there milk in this coffee?
\end{sentenceList}
Between ()and () I do put almond milk in his 
coffee and affirm \q{yes} to ().  I feel it 
proper to read ()'s \q{milk} as really meaning 
\q{almond milk}, in light of ().  Actually 
I should be \i{less} inclined to say \q{yes} 
if (maybe as a prank) someone had instead 
put real (cow) in the coffee.  In responding 
to his question I mentally substitute what 
he almost certainly \i{meant} for how
(taken out of context) () would usually 
be interprted.  In this current 
dialog, the \i{milk} concept not only 
includd vegan milks, apparently, but 
\i{excludes} actual milk.
}
\p{It seems as uf when we are dealing with 
illocutionary force we are obliqued to subject 
what we hear to extra interprtation, rather 
than resting only within \q{literal} meanings 
of sentences, conventionally understood.  
This point is worth emphasizing because it complicates 
are attempts to link illocution with propositional 
content.  Suppose grandma asks me to close the 
kitchen window.  Each of these are plausible and 
basically polite responses: 
\begin{sentenceList}\sentenceItem{} It's not open, but there's still some 
cold air coming through the cracks. 
\sentenceItem{} It's not open, but I closed the window in 
the bedroom.
\sentenceItem{} I can't \mdash{} it's stuck.
\end{sentenceList}
In each case I have not fulfilled her request \visavis{} 
its literal meaning, but I \i{have} acted benevolently 
in terms of conversational maxims.  Many linguists 
seem to analyze hedges like \q{could you please} 
as merely dressing over crude commands: we don't 
want to come across as giving people orders, but 
sometimes we do intend to ask pople to do specific 
things.  As a result, we feel obliged to couch the 
request in conversational gestures that signal 
our awareness of how bald commands may lie outside 
the conversational norms.  These ritualistic 
\q{could you please}-like gestures may have 
metalinguistic content, but \mdash{} so the theory 
goes \mdash{} they do not \i{smantically} alter 
the speech-act's directive nature.
}
\p{The problem with this analysis is that sometimes 
directive and \q{inquisitive} dimensions can 
overlap: 
\begin{sentenceList}\sentenceItem{} Do you have almond milk?
\sentenceItem{} Can you get MsNBC on your TV?
\end{sentenceList}
These \i{can} be read as bare directives, and would 
be interprted as such if the hearer believed the 
speaker already knew that yes, he has almond milk, and yes, 
he gets MsNBC.  They can also be read as bare 
questions with no implicature: maybe fans of 
almond milk and MsNBC endorsing those selections.  
They can \i{also} be read as a mixture of the 
two, as if people expressed themselves like this:
\begin{sentenceList}\sentenceItem{} I think the window is open, can you close it?
\sentenceItem{} I see you have almond milk, can I have some?
\sentenceItem{} If you get MsNBC, can you turn on Rachl Maddow?
\end{sentenceList}
}
\p{I think the mixed case is the most prototypical and pure 
directives or inquiries should be treated as degenerate 
structures whre either directive or inquisitive content 
has dropped out.  After all, even a dictatorial 
command includs the implicit assumption that the order 
both makes sense and is not impossible.  On 
the other hand, we don't ask questions for no 
reason: \q{do you hav almond milk} may be a 
suggestion rather than a request, but it still 
carries an implicature (e.g., that the addressee 
\i{should} get almond milk).   
}
\p{Ordinary requests carry the assumption that addressees 
can follow through without undue inconvenience, 
which includs a package of assumptions about borh 
what is currently the case and what is possible.  
\q{Close the window} only has literal force if the 
window is open.  So when making a request speakers 
have to signal that they recognize the request involves 
certain assumptions and are rational enough to 
accept modifications of these assumptions in 
lieu of literal compliance.  This is why 
interrogative forms liie \q{can you} or 
\q{could you} are both semantically nontrivial 
and metadiscursively polite: they leave open the 
possibility of subsequent discourse framing the original 
request just as a belief-assertion.  Developments 
like \q{can you open the window} \mdash{} \q{no, it's closed} 
are not ruled out.  At the same time, interrogative forms 
connote that the speaker assumes the addressees can 
fulfill the request without great effort: an implicit 
assumption is that they \i{can} and also \i{are 
willing to} satisfy the directive.  This is an 
assumption, not a prsumption: the speaker 
would seem like a bully if he acted as if he 
gave no thought to hid demands being too much 
of and imposition.  This is another rason why requests 
should be framed as question.
}
\p{Sometimes the link between directives and 
belief assertions is made explicit.  A common 
pattern is to use \q{I believe that} as anb
implicature analogous to interrogatives: 
\begin{sentenceList}\sentenceItem{} I believe you have a reservation for Jones?
\sentenceItem{} I believe this is the customer service desk?
\sentenceItem{} I believe we ordered a second basket of garlic bread?
\sentenceItem{} I believe you can help me find a find computer 
acessories in this section?
\end{sentenceList}
These speakers are indirectly signaling what they want 
someone to do by openly stating the requisite 
assumptions \mdash{} \i{I believe you can} in place 
of \i{can you?}.  The implication is that 
such assumptions translate clearly to a 
subsequent course of action \mdash{} the guest who 
\i{does} have that rservation should be checked in; 
the cashier who \i{can} help a customer find 
accessories should do so.  But underlying these 
performances is recognition that 
illocutionary force is tied to background 
assumptions, and conversants are reacting to 
the propositional content of thos assumptions 
as well as the force itself.  If I \i{do} close the 
window I am not only fulfilling 
the request but also confirming that the window 
\i{could} be closed (a piece of information 
that may become relevant in the future).  
}
\p{In sum, when we engage pragmatically with other 
language-users, we tend to do so coopratively, 
sensitive to what they wish to achieve with 
language as well as to the propositional 
details of their discourse.  But this often means 
that I have to interpret propositional 
content in light of contexts and implicatures.  
Note that both of these are possible:
\begin{sentenceList}\sentenceItem{} Do you have any milk?
\sentenceItem{} Yes, we have almond milk.
\sentenceItem{} No, we have almond milk.
\sentenceItem{} Yes \mdash{} actually, we have almond milk.
\sentenceItem{} No, we only have almond milk.
\end{sentenceList}
A request for milk in a vegan restaurant could plausibly be 
interpreted as a request for a vegan milk-substitute.  
So the concept \q{milk} in that context may actually be 
interpreted as the concept \q{vegan milk}.  As Luo 
points out in [], particular concept-maps are 
admissible as in force in specific situations even 
if they deviate noticably from typical usage 
(Luo does not talk about concept \q{maps} but 
about subtyping and various inter-type relations, 
yielding a type-theory of situations I think 
is relevant Conceptual Role theories of situations).  
In any case, responding to the force of speech-acts 
compels me to treat them as not \i{wholly} 
illourionary \mdash{} as in part statements of 
belief (like ordinary assertions).  One reason I need 
to adopt an epistemic (and not just obligatory) 
attitude to illocutionary acts is that I need to 
clarify what meanings the speaker intends, which 
depends on what roles she is assigning to 
constituent concepts.
}
\p{If a diner asks for milk in a vegan restaurant, 
a waiter may plausbily infer that the customer 
believes the restaurant \i{only has} vegan milk, 
so there is no ned to make that explicit; and/or 
she assumes that everyone in the restaurant will hear 
\q{milk} as \q{vegan milk}.  In other words, the 
waiter infers that \q{vegan milk} for her plays the 
same role as \q{milk} for a non-vegan.  This inference 
is not produced by any speech-act subtleties: a related 
inference would be involved in 
\begin{sentenceList}\sentenceItem{} Is there milk in this coffee?
\sentenceItem{} Yes, almond milk.
\end{sentenceList}
Part of reading propositional content is 
syncing our conceptual schemas with our fellow 
conversants.  But the illocutionary 
dimension of a request like \q{can I have some milk?} 
makes this interpretation especially important, 
because the addressee wants to make a good-faoth effort 
to cooperate with the pragmatic intent of the 
spech-act.  But cooperation requires the 
cooperating parties' conceptual schemas to 
be properly aligned.  I therefore have to 
suspend the illocutionary force of a directive 
temporarily and treat it as locutionary 
statement of belief, interpret its apparent 
conceptual undrpinnings in that mode, and 
then add the illocutionary force back in: if I 
brought \i{real} milk to a vegan customer who 
asked for \q{milk} I would be \i{un}-cooprative.
}
\p{The upshot is that conversational implicatures 
help us contextualize the conceptual negotiations 
that guarantee our grasping the correct 
propositional contents, and vice-versa.  This means 
that propositionality is woven throughout both 
assretive and all other mods of language, but it 
also means that propositional content can be 
indecipherable without a detailed picture 
of the current context (including illocutionary 
content).  The proposional content of, 
say, \q{there is milk in this coffee} has to be 
judged sensitive to contexts like \q{milk} 
meaning \q{vegan milk} \mdash{} and this 
propagates from a direct propositional 
to any propositional attitudes which may 
be directed towards it, including requsts like 
\q{plase put milk in this coffee}.
}
\p{Suppose the grandkids close grandma's bedroom window 
when she asks them to close the kitchen window.  
The propositional content at the core of grandma's 
request is that the kitchen window be closed; the 
content attached to it is an unstated belief that 
this window is open.  Thus, the truth-conditions 
satisfying her implicit understanding would be 
that the kitchen window went from being open to being 
closed.  As it happens, that window is already 
closed.  So the truth-conditions that would satisfy 
grandma's initial belief-state do not obtain \mdash{} her 
beliefs are false \mdash{} but the truth conditions satisfying 
her desired result \i{do} obtain.  The window 
\i{is} closed.  Yet the grand kids should not thereby 
assume that her request has been properly responded to; 
it is more polite to guess at the motivation behind 
the request, e.g., thay she felt a draft 
of cold air.  In short, they should look outside 
the truth conditions of her original 
request taken literally, and \i{interpret} 
her requesting, finding different content 
with different truth-conditions that are both consistent 
with fact and address whatever pragmatic goals 
grandma had when making her request.  They might 
infer her goal is to prevent an uncomfortable 
draft, and so a reasonable \q{substitute content} is 
the proposition that \i{some} window is open, 
and they should close \i{that} one. 
}
\p{So the grandkids should reason as if translating 
between these two implied meanings: 
\begin{sentenceList}\sentenceItem{} I believe the kitchen window 
is open \mdash{} please close it!
\sentenceItem{} I believe some window 
is open \mdash{} please close it!
\end{sentenceList}
They have to revise the simplest reading of 
the implicit propositional content of grandma's 
\i{actual} request, because the actual request is 
inconsistent with facts.  In short they 
feel obliged to explore propositional alternatives 
so as to find an alternative, implicit request whose 
propositional content \i{is} consistent with 
fact and also meets the original request's illocutionary 
force cooperatively.  
}
\p{In essence, we need to express a requester's desire as 
itslf, in its totality, a specific propositional content, 
thinking to purselves (or even saying to others) things 
like
\begin{sentenceList}\sentenceItem{} Grandma wants us to close the window.
\sentenceItem{} He wants a bottle opener.
\end{sentenceList}
But to respond politely we need to modify 
the parse of their requests to capture the 
\q{essential} content:
\begin{sentenceList}\sentenceItem{} Grandma wants us to eliminate the cold draft.
\sentenceItem{} He wants something to open that bottle.
\end{sentenceList}
We have to read outside the literal unterpretation 
of what they ar saying.  This re-reading is something 
that may be appropriate to do with respect to 
other forms of speech also: sometimes 
the true gist of what someone wants to communicate 
is not stated directly:
\begin{sentenceList}\sentenceItem{} I think you could do excellent work in 
this class, and I think you are doing pretty well.
\sentenceItem{} I am not going to talk about the refs because 
I don't want to get fined.
\sentenceItem{} If she wants to win the nomination she 
needs to be as charismatic on the campaign trail as 
she was during the debate.
\end{sentenceList}
But our conversational responsibilty to infer 
some unstated content is especially pronounced 
when we are rsponding to an explicit 
request for something. 
}
\p{Certainly, in any case, meanings are not literal.  
But how then do we understand what people are saying?  
Trying to formulate a not-entirely-haphazard 
account of this process, we can spculate 
that interpreting what someone is \q{really} saying 
involves systematically mapping their eapparent 
concepts and references to some superimposed 
inventory designed to mitigate false beliefs or 
conceptual misalignments among language users in some 
context.  That means, we are looking for mappings 
like \i{milk} to \i{almond milk} in () from a 
vegan restaurant, or \i{kitchen window} to 
\i{bedroom window} in () if it is the latter 
that is open:
\begin{sentenceList}\sentenceItem{} Can I have some milk?
\sentenceItem{} Can you close the kitchen window?
\end{sentenceList}
The point of thse \q{mappings} is that they 
preserve the possibility of 
modeling the \i{original} propositional content 
by identifying truth conditions 
for that content to be satisfied.
}
\p{A \i{literal} truth-condition model doesn't work in 
cases like () and (): th diner's request 
is \i{not} satisfied if it is the case 
that there is now (real) milk in her coffee, and 
grandma's request is not necessarily satisfied if it is 
the case that the kitchen window is closed.  The 
proposition \q{the kitchen window is closed} only bears on 
grandma's utterance insofar as she believes that 
this window is open and causing a draft.  So if we want 
to interpret the underlying locutionary content 
of () and () truth-theoreticaly we need to 
map the literal concepts appearing 
in thse sentences to an appropriate translation, 
a kind of \q{coordinate transformation} that 
can map concepts onto others, like milk/almond milk 
and kitchen window/bedroom window.  
}
\p{Simultaneous with propositional content, of 
course, are attitudes: the difference between 
asserting and wanting that the window is closed.  
It is hard to deny that \i{some} propositional 
content is involved with each linguistic epxression, 
because simply by being a structured mental activity 
the effort to formulate sentencs must be extended with 
some purpose.  We say (and write) things to help make 
something or other the case.  But there are several 
challenges to disentangle the role that propositional 
content actually plays in meaning.  One problem I 
just considered is that the right propositional content 
does not always come from \i{literal} meaning: the 
vegan \i{doesn't} want real milk in her coffee.  
The idea of \q{mapping} is one away to address this: 
in place of \q{literal} meaning we can substitute 
meanings undr \q{coordinate} transforms, where 
concepts transition from their literal designation 
to their roles: the vegan wants the product that 
plays the \i{conceptual role} of \q{milk} in her 
own frame of reference (at least in the context of, 
say, dining, as opposed to a context like 
checking whether a pit bull is lactating).  
But there are two other concerns we should have 
about propositional content, which 
I will discuss to close out this section.
}
\subsection{The Problem of Opaque Truth Conditions}
\p{My analysis related to conceptual \q{transforms} 
assumed that we can find, substituting for \i{literal} 
propositional content, some \i{other} 
(reprsentation of a) proposition that fulfills a 
speaker's unstated \q{ral} meaning.  Sometimes 
this makes sense: the proposition that the 
\i{bedroom} window is closed can neatly, 
if the facts warrant, play the role of the 
proposition that the kitchen window is closed.  
But we can run the example differently: there 
may be \i{no} window open, but instead a draft 
caused by non-airtight windows (grandma might ask 
us to put towels by the cracks).  Maybe there is 
no draft at all (if grandma is cold, we can 
fetch her a sweater).  Instead of a single 
transform, we need a a system of potential transforms 
that can adapt to the facts as we discover them  
Pragmatically, the underlying problem is 
that grandma is cold.  We can address this 
\mdash{} if we want to faithfully respond to her request, 
playing the role of cooperative conversation 
partners (and grandkids) \mdash{} via a matrix 
of logical possibilities: 
\begin{sentenceList}\sentenceItem{} If the kitchen window is closed, 
we can see if other windows are open.
\sentenceItem{} If no windows are open, 
we can see if there is a draft through the window-cracks.
\sentenceItem{} If there is no draft, we can 
ask if she wants a sweater.
\end{sentenceList}
This is still a logical process: starting from an 
acknowledged proposition (grandma is cold) we 
entertain various other propositional possibilities, 
trying to rationally determine that pragmas we 
should enact to alter that case 
(viz., to instead make true the proposition that 
grandma is warm).  Here we are not just testing 
possibilities against fact, but strategically 
acting to modify some facts in our environment.  
}
\p{The kind of reasoning involved here is not logical 
reasoning per se: abstract logic does not tell us 
to check th bdroom window if the kitchen window is closed, 
or to check for gaps and cracks if all windows 
are closed.  That is practical, domain-specific knowledge 
about windows, air, weather, and houses.  
But we are still deploying our practical 
knowledge in logical ways.  There is a logical 
structure underpinning grandma's request and 
our response to it.  In sum: we (the grandkids) 
are equipped with some practical knowldge 
about houses and a faculty to logically utilize this 
knowledge to solve the stated problem, reading 
beyond the \i{explicit} form of grandma's discourse.  
We use a combination of logic and background 
knowledge to reinterpret the discourse as needed.  
By making a request, grandma is not expressing 
one attitude to one proposition, so much as 
\i{initiating a process}.  This is why it would 
be impolite to simply do no more 
if the kitchen window is closed: our conversational 
responsibility is to enact a process trying to 
redress grandma's discomfort, not to entertain the 
truth of any one proposition.
}
\p{For all that, there is still an overarching logical 
structure here that languages clearly marshals.  
But once we converge on the \q{language initiates 
a process} model, we can find examples where the 
logical scaffolding gets more tenuous.  
Consider:
\begin{sentenceList}\sentenceItem{} My colleague Ms. O'Shea would like to interview 
Mr. Jones, who's an old friend of mine.  Can he take this call?
\sentenceItem{} I'm sorry, this is his secretary.  Mr. Jones 
is not available at the moment.
\end{sentenceList}
It sounds like Ms. O'Shea is trying to use personal 
connections to score an interview with Mr. Jones.  Hence 
her colleague initiates a process intended to 
culminate in Ms. O'Shea getting on the telephone 
with Mr. Jones.  But his secretary demurs with a 
familiar phrase, deliberately formulated to 
foment ambiguity: () could mean that Mr. Jones 
is not in the office, or in a meeting, or 
unwilling to talk, or even missing (like 
the ex-governer consummating an affair in Argentina 
while his aides thought he was hiking in Virginia).  
Or:
\begin{sentenceList}\sentenceItem{} Mr. Jones, were you present at a meeting where 
the governer promised your employer 
a contract in exchange for campaign contributions?
\sentenceItem{} After consulting with my lawyerrs, I decilne 
to answer that question on the grounds that it 
may incriminate me.
\end{sentenceList} 
Here Mr. Jones neither confirms nor denies his 
presence at a corrupt meeting.
}
\p{As these examples intimate, the processes language 
initiates do not always result in a meaningful 
logical structure.  But this is not necessarily 
a complete breakdown of language:
\begin{sentenceList}\sentenceItem{} Is Jones there?
\sentenceItem{} He is not available.
\end{sentenceList}
The speaker of () does not 
provide any logical content: it neither affirms 
nor denies Jones's presence.  Nonetheless that speaker 
is a cooperative conversational partner 
even if they are not being cooperative in real life): 
() responds to the implicature in () that what the 
first speaker really wants is (for instance) 
to interview Jones.  
So the second speaker conducts what I called 
a \q{transform} and maps \q{Jones is here} to 
\q{Jones is willing to be interviewed}.  
Responding to this \q{transformed} question allows 
() to be (at least) linguistically cooperative 
while nonetheless avoiding a response at the 
\i{logical} level to ().  () obeys 
conversational maxims but is still rather obtuse.
}
\p{This pattern of logical evasiveness might seem to be 
endemic only to slippery human languag, but analogous 
exsamples can be found even in the strict milieu of 
computer programming.  Suppose 
I am writing a function which counts the 
number of non-blank, non-comment 
lines in a file.  My implementation might start like this 
(in pseudo-code):
\begin{sentenceList}\sentenceItem{} File f = File.open(path, File.READ\_ONLY);
\sentenceItem{} if f.isEmpty() return 0;
\end{sentenceList}
So the first line tries to open a file with the 
specified path, and the second line checks if the file 
is indeed open and non-empty.  If not, it returns the 
number zero, meaning there are no non-blank non-comment 
lines in the file.
}
\p{In a typical application framework a file will be 
opened if it exists and if the current user has permission 
to read and/or write to/from the file (here the necessary 
permission is read-only).  For sake of presentation 
I assume that a file is considered non-empty 
only if it is open and has some content (i.e., you 
can read something from the file).  So if 
the file \i{is} empty, that can mean several things: 
either it does not exist; or it cannot be opened because 
of inadequate permissions; or it \i{can} be 
opened but has no content.  My function is 
noncommital and returns zero in each of these cases.  
In particular, I gloss over someinformation: 
the zero return value is analogous to Mr. Jones 
being \q{unavailable} for an interview.  
}
\p{In the formal computer setting, we are not allowed to 
logicaly infer what expressions \q{really} mean: we 
should not, for instance, if some file does not 
exist, instead try to opn a different file 
with a similar name.  Therefore the \q{meaning} of 
code expressions should be tied explicitly to 
individual logical conditions.  For cample, the 
\q{meaning} of line () should be tied to the proposition 
that the file (whose location is specified by \i{path}) 
is open.  However, the code never actually enagages with 
that proposition.  I do not necessarily determine whether 
the file is open, can be opened, or even whether it 
exists. 
}
\p{On this basis, it seems as if the \q{meaning} of my 
call to open the function is not a matter of 
ascertaining or bringing about certain truth 
conditions.  It is true that my instruction 
may bring about certain truths (specifically, make it true 
that the file is open).  But we should 
not conclude that this potential state of affairs is 
what the instruction \i{means}.  As the programmer, I 
do not \q{want} the file to be opened; I have no 
vested interest (this is not like my wanting milk 
for coffee or to open a beer bottle).  Instead, 
I am an intermediary between application-users and 
the system kernel: my role is translate what \i{users} 
want into system instructions.  Granted, presumably the 
\i{user} wants the file she's interested in to be opened 
(as an intermediate step toward getting information 
from that file).  But what I contribite as 
code are \i{instructions}, and instructions have 
an effect on the state of the overall computational 
environment where my application is hosted.  
In the general case I am not aware of eactly what 
new state obtains.  Attempting to open a file 
may cause it to be opened, but it may also 
cause the software reprsentation of that 
file \mdash{} a so-called \q{handle} \mdash{} to acquire a flag 
indicating that the referenced file does not exist 
(i.e., its path descrribes a nonexistent location), 
or that it does exist but cannot be opened due to 
insufficient permission, or that the permissions 
are satisfactory but there is a temporary lock 
from another application writing to the file.  
If needed I could attempt to ascertain the file's state 
at this level of detail, but it turns out to be 
irrelevant to my own algorithm.  In short, 
insofar as the \q{meaning} of computer code are the 
instructions it emits, these meanings correspond 
to state-changes only in coarse-grained ways.  
There may be propositional content associated to each possible 
state \mdash{} there are propositions that 
the file is open, or nonxistent, or inaccessible, 
or temporarily locked \mdash{} but the code does not engage 
the question of \i{which} proposition is 
(or becomes) true.     
}
\p{Underlying both these formal and natural 
language examples is the idea that the meanings 
of expressions are associated with changes of 
something's state: computer statements are instructions 
to change state and a large class of linguistic 
expressions are requests to do so.  Whenever 
there is a state-change there is a corresponding 
proposition to the effect that the \q{thing} is 
now in state \Stwo{} whereas it was before in \Sone{}.  
Even if the instruction/request cannot be fulfilled there 
is the concordant proposition that the thing is 
still in \Sone{} can can \i{not} be brought into 
\Stwo{}.  So it is trivial to read logical structures 
onto linguistic evetualities, by leveraging 
the idea that for any conceivable 
state of any conceivable state-bearer there exists a 
proposition that such bearer is in such state.  
}
\p{The problem for truth-theoretic semantics, as 
I see it, is that these trivial state-to-proposition 
conversions are just that \mdash{} theoretically trivial.  
We should not care about a \i{trivial} truth-semantic 
theory.  If we have a semantic theory wherein, let's 
say, \q{meanings} are really \i{initiators of 
state-change processes}, then we can trivially 
convert this into a truth-theoretic theory.  But 
that is not an intersting truth-theoretic semantics; 
it is a trivial truth-theoretic theory grafted onto 
an interesting \q{state change} theory.  
}
\p{So much is not to call truth-theoretic semantics 
uninteresting.  But for us to take truth-theoretic semantics 
seriously we need to accept the idea that this 
paradigm can \i{motivate} analysis qhich leads to 
interesting results, taking us somewhere we may not 
arrive otherwise.  The fact that a given semantic theory 
has some formulaic translation into truth-thoretic terms 
does not guarantee that truth-theoretic intuitions 
actually play an important role in that other theory.  
}
\p{For truth-theoretic intuitions to be legitimately 
consequential toward a semantic theory, we need to 
ascertain to what degree logical structures 
actually play a cognitive role in how we use 
language to accomplish things in the world.  Obviously, 
as rational beings our thought processes will 
be informed by logic and to some degree can be 
retroactively modeled via logical complexes.  
But \q{logic} appears to play a role in these 
cognitive operations only inidrectly.  There seems 
to be some medium \mdash{} perhaps conceptual roles, 
or state-changes \mdash{} that \q{carries} logic 
into the cognitive realm.  We should reject   
truth-theoretic semantics if it seems to 
proceed as if that \q{medium} can be sidetracked 
\mdash{} that we can analyze a logical form in language 
directly, without analyzing the vehicle 
by which logical considerations can enter 
language processing.
}
\subsection{Truth Conditions are not Polar}
\p{My second quibble with truth-theoretic semantics 
is that it relies on a certain \faconaparler{}.  
We (in the contxt of analysis, not at her house) might 
say something like:
\begin{sentenceList}\sentenceItem{} Grandma wants the proposition \q{the kitchen 
window is closed} to be true.
\sentenceItem{} Grandma wants the proposition 
\q{I am cold} to be false.
\end{sentenceList}
My prior analysis focusd on the fact that the more relevant 
proposition is that in (), and we have to 
read () from ().  But it should also be obvious 
that grandma does not care about \i{propositions}.  
She only cares about \q{I am cold} as a \i{proposition} 
because she does not want to be cold.
}
\p{By contrast, sometimes people care about propositions 
\i{as proposition}.  A mathematician who has 
staked her reputatuon on a conjecture may want 
the conjuecture to be true.  But this desirre is not 
like the desire of a Sanders or a Toronto Maple 
Leafs supporter to want propositions like 
\q{Sanders has won} or \q{Toronto has won} to 
b true.  The mathematician's conjectur may 
be in an obscure field where there are no 
apparent real-world consequences with 
benefits that can proceed from truth rather than 
falsehoodb (it's not like, say, the conjecture 
will allow her to prove the validity of an 
encryption sheme she can monetize).
}
\p{One way to put this is that the mathematician 
desires a \i{proposition} to be true, whereas the 
Leafs or Sanders supporters 
want certin \i{propostional content} to be true.  
But if wechave to bring in propositional 
content, our analysis \mdash{} so it 
sems \mdash{} becomes circular.  If the propositional 
content of \q{the Leafs win the cup} is that th Leafs 
win th cup, then a fans who wants the Leafs to 
win the cup obviously wants the propositional content 
of \q{the Leafs win the cup} to be true \mdash{} but that's 
becaus the idea of the Leafs winning the cup 
\i{is} the propositional content of 
\q{the Leafs win the cup}. 
}
\p{Let's say that there are two inverse operators: 
one maps an idea onto a proposition, 
and one maps the proposition back onto the 
idea, which is its \q{content}.  A semantic 
analysis would be circular if it just immediatly 
reversed one operator with the other.  Having said that, 
the reversal may be \i{separated} by several steps, making 
it non-trivial.  Suppose the Leafs ar playing 
the Jets:
\begin{sentenceList}\sentenceItem{} Are you rooting for the Jets?
\sentenceItem{} Well, I want the Leafs to win.
\end{sentenceList}
The second speaker indirecly answers the first in the negative.  
This is conversationally reasonable \mdash{} it does not violate any 
Naxim of Relevance \mdash{} because the first 
speaker has available a chain of inference like:
\begin{sentenceList}\sentenceItem{} There is a proposition (call it \POne{}) 
whose propositional content is that the Jets win 
and a proposition (\PTwo{}) whose content is that the Leafs win.
\sentenceItem{} \PTwo{} entails the negation of \POne{}.
\sentenceItem{} He wants the content of \PTwo{} to be true.
\sentenceItem{} His wishes are only consistent 
with \POne{} being false.
\sentenceItem{} He wants the content of \POne{} to be false.
\end{sentenceList}
The key step here is at () and (): in these steps 
of reasoning the predicate relation between propositions 
is expressly thematized.  In the course of 
understanding language, there may be occasions when we 
need to identify logical connectives as implicit 
to the conceptual framework which a speaker is obviously 
assuming.  In that case we are working with 
propositions rather than propositional content: 
there is nothing in the idea of a Leafs victory 
\i{apriori} that contains the idea of a Jets 
defeat.  The negative entailment relation only arises 
in particular situations \mdash{} when the two teams 
play each other (and implicitly further restrictions, 
like there are no draws or suspended games). 
}
\p{It is, in short, a \i{feature} of a specific situation that 
the idea \q{Leafs win} entails the negation of 
\q{Jets win}; ergo, logical connectives are 
one facet of situational models \i{sometimes} 
relevant to linguistic processes.  Similarly, speaker 
sentiment is \i{sometimes} relevant.  But there 
is a structural isolation between these \q{systems}: 
those \q{processing units} that can bear speaker 
sentiment (polarity) are different 
from those that can bear logical connectives 
(except in unusual cases, like the mathematician 
and the conjecture).
}
\p{When we analyze a fan's sentiment \i{wanting} the 
Leafs to win, we are analyzing polarity \visavis{} 
propositional \i{content}.  When we observe  
negative entailment, we are analyzing relations 
among \i{propositions}.  The proposition 
\q{encapsulates} the content so it can be part 
of logistic structures, where connectives like 
entailment and disjunction make sense.  In 
computer science, the technical pattern of 
such \q{encapsulation} is often 
called a \q{monad}, and monadic analysis has been 
adopted in linguistics as well, following 
Chung-Chieh Shan.  We can say that logical propositions 
\i{per se} are monadic packages that \q{wrap} propositional 
content, subject it to some logical manipulation, 
and at some later point in processing \q{retrieve} 
the content from the proposition.  This is not circular 
because the content is not \i{immediately} extracted from 
th \q{monad}.  In (), the content wrapped in 
\POne{} is held through several processing steps before yielding 
the final interpretation (he wants the Jets to lose).  
The final extraction corresponds to transforming 
\q{He wants the proposition that the Jets lose to be false} 
to \q{He wants the Jets to lose}.  Stated side-by-side, 
this transform is redunant.  The difference is that 
here the original proposition \Ptwo{} is not 
directly asserted as a linguistic meaning; it rather  
falls out of a logical process.  
}
\p{On this analysis, there is a \i{part}, or 
\i{substructure}, of linguistic processing that involves 
wrapping propositional content into propositions.  
But these wrappings are only meaningful 
when there is some \q{delay} between procssing 
steps \mdash{} essentially, when there is some 
explicit sense that a given idea stands in some 
well-defined logical reoation to some other idea.  
Computer programers can crate \q{trivial} monads 
whose behaviors do not deviate at all from an 
imperative style of programming, but any code 
written with \q{trivial} monads can be refactored 
such that the monads disappear.  This is analogous 
to the circularity between propositional 
content and propositions: the \i{meaning} of 
a proposition \i{is} its content, so 
the proposition \q{monad} has processing significance 
only when the meaning itself needs to be \q{held} 
awaiting the resolution of some logical nexus.  
If there is no processing structure that demands 
the content to be \q{held}, then the proposition 
just \q{decays} to its content, and 
essentially disappears. 
}
\p{These points suggest that while modeling linguistic 
meanings in terms of propositions is \i{sometimes} 
appropriate, in the egneral case it is merely 
circular or tautological.  We certainly should not 
put forth a theory that the \q{meaning} of an 
idea is its proposition \mdash{} the turth is more 
the opposite.  The meaning of a proposition is 
its propositional content.  So, if we want a 
truth-theoretic semantics that is 
applicable for general cases \mdash{} not just especially 
logically ordered situations, like a winner-take-all 
sporting match \mdash{} we need a truth-conditional 
theory of propositional \i{content} separate and apart 
from a theory of propositions.
}
\p{I am discinclined to believe that such a theory is possible, 
since I think we can give a theory of propositional 
content but I don't think logic would 
figure strongly in it.  Having said that, I should 
now explain what such a theory of propositional 
content \i{should} look like, and I'll lave it as a 
(mostly) rhetorical question whether this theory 
does or does not leverage logic or something else.
}

\end{document}
